% Options for packages loaded elsewhere
\PassOptionsToPackage{unicode}{hyperref}
\PassOptionsToPackage{hyphens}{url}
\documentclass[
]{article}
\usepackage{xcolor}
\usepackage{amsmath,amssymb}
\setcounter{secnumdepth}{-\maxdimen} % remove section numbering
\usepackage{iftex}
\ifPDFTeX
  \usepackage[T1]{fontenc}
  \usepackage[utf8]{inputenc}
  \usepackage{textcomp} % provide euro and other symbols
\else % if luatex or xetex
  \usepackage{unicode-math} % this also loads fontspec
  \defaultfontfeatures{Scale=MatchLowercase}
  \defaultfontfeatures[\rmfamily]{Ligatures=TeX,Scale=1}
\fi
\usepackage{lmodern}
\ifPDFTeX\else
  % xetex/luatex font selection
\fi
% Use upquote if available, for straight quotes in verbatim environments
\IfFileExists{upquote.sty}{\usepackage{upquote}}{}
\IfFileExists{microtype.sty}{% use microtype if available
  \usepackage[]{microtype}
  \UseMicrotypeSet[protrusion]{basicmath} % disable protrusion for tt fonts
}{}
\makeatletter
\@ifundefined{KOMAClassName}{% if non-KOMA class
  \IfFileExists{parskip.sty}{%
    \usepackage{parskip}
  }{% else
    \setlength{\parindent}{0pt}
    \setlength{\parskip}{6pt plus 2pt minus 1pt}}
}{% if KOMA class
  \KOMAoptions{parskip=half}}
\makeatother
\usepackage{longtable,booktabs,array}
\newcounter{none} % for unnumbered tables
\usepackage{calc} % for calculating minipage widths
% Correct order of tables after \paragraph or \subparagraph
\usepackage{etoolbox}
\makeatletter
\patchcmd\longtable{\par}{\if@noskipsec\mbox{}\fi\par}{}{}
\makeatother
% Allow footnotes in longtable head/foot
\IfFileExists{footnotehyper.sty}{\usepackage{footnotehyper}}{\usepackage{footnote}}
\makesavenoteenv{longtable}
\usepackage{graphicx}
\makeatletter
\newsavebox\pandoc@box
\newcommand*\pandocbounded[1]{% scales image to fit in text height/width
  \sbox\pandoc@box{#1}%
  \Gscale@div\@tempa{\textheight}{\dimexpr\ht\pandoc@box+\dp\pandoc@box\relax}%
  \Gscale@div\@tempb{\linewidth}{\wd\pandoc@box}%
  \ifdim\@tempb\p@<\@tempa\p@\let\@tempa\@tempb\fi% select the smaller of both
  \ifdim\@tempa\p@<\p@\scalebox{\@tempa}{\usebox\pandoc@box}%
  \else\usebox{\pandoc@box}%
  \fi%
}
% Set default figure placement to htbp
\def\fps@figure{htbp}
\makeatother
\ifLuaTeX
  \usepackage{luacolor}
  \usepackage[soul]{lua-ul}
\else
  \usepackage{soul}
\fi
\setlength{\emergencystretch}{3em} % prevent overfull lines
\providecommand{\tightlist}{%
  \setlength{\itemsep}{0pt}\setlength{\parskip}{0pt}}
\usepackage{bookmark}
\IfFileExists{xurl.sty}{\usepackage{xurl}}{} % add URL line breaks if available
\urlstyle{same}
\hypersetup{
  hidelinks,
  pdfcreator={LaTeX via pandoc}}

\author{}
\date{}

\begin{document}

\pandocbounded{
\includegraphics[keepaspectratio]{./media/image2.jpeg}}

o

\begin{quote}
\textbf{TRƯỜNG ĐẠI HỌC NÔNG LÂM TP HCM}

\textbf{KHOA CÔNG NGHỆ THÔNG TIN}
\end{quote}

TIỂU LUẬN TỐT NGHIỆP

\begin{quote}
\textbf{TÊN ĐỀ TÀI}
\end{quote}

\textbf{Nghiên cứu cải tiến TickNets bằng cơ chế}

\textbf{attention CBAM cho bài toán phân loại ảnh}

\begin{quote}
\textbf{Ngành : Công nghệ thông tin.}

\textbf{Niên khoá : 2022 -- 2026.}

\textbf{Lớp : DH22DTC.}

\textbf{Sinh viên thực hiện : Chu Toàn Đức.}
\end{quote}

\pandocbounded{
\includegraphics[keepaspectratio]{./media/image2.jpeg}}

\textbf{MẪU BÁO CÁO LUẬN VĂN, TIỂU LUẬN TỐT NGHIÊ}

\begin{quote}
\textbf{TRƯỜNG ĐẠI HỌC NÔNG LÂM TP HCM}

\textbf{KHOA CÔNG NGHỆ THÔNG TIN}

TIỂU LUẬN TỐT NGHIỆP

\textbf{TÊN ĐỀ TÀI}
\end{quote}

\textbf{Nghiên cứu cải tiến TickNets bằng cơ chế}

\textbf{attention CBAM cho bài toán phân loại ảnh}

\subsection{CÁN BỘ HƯỚNG DẪN SINH VIÊN THỰC
HIỆN}\label{cuxe1n-bux1ed9-hux1b0ux1edbng-dux1eabn-sinh-viuxean-thux1ef1c-hiux1ec7n}

\begin{quote}
TS. Nguyễn Văn Dũ 1. Chu Toàn Đức (22130047).

TP.HỒ CHÍ MINH,24 tháng 06 năm 2026

\textbf{DANH SÁCH CHỮ VIẾT TẮT}

NLP \textbf{N}atural \textbf{L}anguage \textbf{P}rocessing Xử lý ngôn
ngữ tự nhiên.

CNN \textbf{C}onvolutional \textbf{N}eural \textbf{N}etwork Mạng nơ-ron
tích chập.
\end{quote}

CV \textbf{C}omputer \textbf{V}ision.

Thị giác máy tính.

\begin{quote}
DL \textbf{D}eep \textbf{L}earning Học sâu.

ANN \textbf{A}rtificial \textbf{N}eural \textbf{N}etwork Mạng nơ-ron
nhân tạo.

FC \textbf{F}ully \textbf{C}onnected Layer Lớp đầy đủ.

CAM \textbf{C}hannel \textbf{A}ttention \textbf{M}odule Mô-đun chú ý
theo kênh.

SAM \textbf{S}patial \textbf{A}ttention \textbf{M}odule

Mô-đun chú ý theo không gian.

KNN \textbf{K}-\textbf{N}earest \textbf{N}eighbors K láng giềng gần
nhất.
\end{quote}

\section{DANH MỤC HÌNH ẢNH}\label{danh-mux1ee5c-huxecnh-ux1ea3nh}

\begin{quote}
Hình 2.1. Minh họa về Neural Network 5

Hình 2.2 Ví dụ về kiến trúc của CNN 6

Hình 2.3 Ví dụ về Output của Convolution 8

Hình 2.4 Ví dụ về việc sử dụng hàm kích hoạt ReLU 9

Hình 2.4 Ví dụ về Max-pooling và Average-pooling 11

Hình 2.5 Minh họa Fully Connected Layer 12

Hình 2.6 Kiến trúc mô hình VGG16 13

Hình 2.7 Ví dụ cách thức hoạt động của CNN 14

Hình 2.8 Công thức tính toán output của Convolution 15

Hình 2.9 Mô hình hoạt động của CAM 17

Hình 2.10 Mô hình hoạt động của SAM 19

Hình 2.11 Mô hình CBAM 20

Hình 3.1a Tổng quan bài toán phân loại ảnh

Hình 3.1b Biểu diễn hình ảnh trong máy tính

Hình 3.1c Bốn thách thức chính của bài toán phân loại ảnh

Hình 3.1d Các lĩnh vực ứng dụng chính của phân loại hình ảnh trong thực tiễn

Hình 3.2 Mô hình CNN kết hợp CBAM 23

Hình 3.3 Quá trình xây dựng và huấn luyện mô hình 24

Hình 3.4 Tải tập dữ liệu 26

Hình 3.5 Mô hình CNN cho tập dữ liệu MNIST-Fashion 28

Hình 3.6 Mô hình CNN cho tập dữ liệu Cifar-10 29

Hình 3.7 Cấu trúc Confusion matrix 31

Hình 3.8 Biểu đồ loss và accuracy của mô hình CNN có áp dụng Attention
và không áp dụng với tập dữ liệu MNIST-Fashion 33

Hình 3.9 Biểu đồ loss và accuracy của mô hình CNN có áp dụng Attention
và không áp dụng với tập dữ liệu Cifar-10 34
\end{quote}

\section{DANH MỤC BẢNG}\label{danh-mux1ee5c-bux1ea3ng}

\begin{quote}
Bảng 3.1 Bảng mô tả dữ liệu tập dữ liệu MNIST-Fashion 25

Bảng 3.2 Bảng mô tả dữ liệu tập dữ liệu Cifar-10 26

Bảng 3.3 Bảng so sánh kết quả giữa các giải thuật trên tập dữ liệu
MNIST-Fashion

. 35

Bảng 3.4 Bảng so sánh kết quả giữa các giải thuật trên tập dữ liệu
Cifar-10 35

\textbf{TÓM TẮT}

Bài báo cáo này nghiên cứu, đề xuất và hiện thực hóa giải pháp cải tiến
kiến trúc mạng nơ-ron tích chập nhẹ TickNets bằng cách tích hợp phân cấp
cơ chế chú ý khối tích chập CBAM (Convolutional Block Attention Module)
cho bài toán phân loại hình ảnh. Thông qua các thử nghiệm thực tiễn trên
hai nhóm tập dữ liệu chuẩn học thuật Fashion-MNIST, CIFAR-10 và ứng dụng
thực tế phức tạp PlantVillage, Chest-Xray-Pneumonia, kết quả thực nghiệm
chỉ ra rằng giải pháp chèn CBAM phân cấp tại các điểm nối mạng không chỉ
giúp loại bỏ hiệu quả nhiễu nền trên ảnh bối cảnh thực tế mà còn khắc
phục tối ưu sự suy giảm hiệu năng cục bộ trên ảnh kích thước nhỏ. Dựa
trên các độ đo chuẩn mực bao gồm accuracy, precision, recall và
f1-score, mô hình cải tiến đề xuất mang lại hiệu suất vượt trội và ổn
định hơn đáng kể so với các mô hình baseline không áp dụng hoặc chèn cơ
chế chú ý sai vị trí, đóng góp một hướng đi có giá trị khoa học cao cho
các ứng dụng học sâu tối giản trên thiết bị có tài nguyên hạn chế.
\end{quote}

\section{MỤC LỤC}\label{mux1ee5c-lux1ee5c}

\begin{quote}
DANH SÁCH CHỮ VIẾT TẮT

DANH MỤC HÌNH ẢNH

DANH MỤC BẢNG

TÓM TẮT

MỤC LỤC

MỞ ĐẦU
\end{quote}

\begin{enumerate}
\def\labelenumi{\arabic{enumi}.}
\item
  LÝ DO CHỌN ĐỀ TÀI
\item
  MỤC TIÊU VÀ PHẠM VI NGHIÊN CỨU
\item
  Ý NGHĨA KHOA HỌC VÀ THỰC TIỄN CỦA ĐỀ TÀI
\end{enumerate}

\begin{quote}
CHƯƠNG 1: TỔNG QUAN ĐỀ TÀI
\end{quote}

\begin{quote}
CHƯƠNG 2: CƠ SỞ LÝ THUYẾT
\end{quote}

\begin{enumerate}
\def\labelenumi{2.\arabic{enumi}.}
\item
  MẠNG NƠ-RON TÍCH CHẬP SIÊU NHẸ TICKNETS
  \begin{enumerate}
  \def\labelenumii{2.1.\arabic{enumii}.}
  \item
    Khái niệm và đặc điểm của họ mạng tích chập siêu nhẹ TickNets
  \item
    Các phép tích chập tối giản (Depthwise Convolution và Pointwise Convolution)
  \item
    Kiến trúc khối Full-Residual Point-Depth-Point (FR-PDP)
  \item
    Cấu trúc xương sống co giãn kênh hình dấu tích (Tick-shape Backbone) và mô hình TickNet-small
  \item
    Cơ chế lan truyền đặc trưng và cách thức hoạt động của TickNets
  \end{enumerate}
\item
  CƠ CHẾ CHÚ Ý (ATTENTION MECHANISM)
  \begin{enumerate}
  \def\labelenumii{2.2.\arabic{enumii}.}
  \item
    Cơ chế chú ý theo kênh Squeeze-and-Excitation (SE) mặc định trong TickNets
  \item
    Mô-đun chú ý theo kênh (Channel Attention Module - CAM) của CBAM
  \item
    Mô-đun chú ý theo không gian (Spatial Attention Module - SAM) của CBAM
  \item
    Mô-đun chú ý khối tích chập CBAM hoạt động tuần tự (CAM + SAM)
  \end{enumerate}
\item
  PHƯƠNG PHÁP NGHIÊN CỨU
\end{enumerate}

\begin{quote}
CHƯƠNG 3: BÀI TOÁN PHÂN LOẠI HÌNH ẢNH VÀ ĐỀ XUẤT CẢI TIẾN PHÂN CẤP
\end{quote}

\begin{enumerate}
\def\labelenumi{3.\arabic{enumi}.}
\item
  GIỚI THIỆU BÀI TOÁN PHÂN LOẠI HÌNH ẢNH
\item
  ĐỀ XUẤT MÔ HÌNH TICKNETS KẾT HỢP CƠ CHẾ CHÚ Ý PHÂN CẤP (HIERARCHICAL ATTENTION)
  \begin{enumerate}
  \def\labelenumii{3.2.\arabic{enumii}.}
  \item
    Thiết kế kiến trúc đề xuất TickNet-small CBAM-Hook
  \item
    Cơ chế chú ý phân cấp (Giữ SE nội khối FR-PDP, chèn CBAM ngoại khối tại các điểm nối Hooking points và cuối mạng)
  \item
    Phân tích toán học giải quyết sự chồng chéo đặc trưng không gian giữa Depthwise Convolution và Spatial Attention Module (SAM)
  \end{enumerate}
\item
  HIỆN THỰC GIẢI PHÁP VÀ THỰC NGHIỆM
  \begin{enumerate}
  \def\labelenumii{3.3.\arabic{enumii}.}
  \item
    Môi trường thực nghiệm và xây dựng mô hình
    \begin{enumerate}
    \def\labelenumiii{3.3.1.\arabic{enumiii}.}
    \item
      Môi trường thực nghiệm (Nền tảng Kaggle, GPU T4, TensorFlow và Keras)
    \item
      Mô tả 4 bộ dữ liệu thực nghiệm
      \begin{enumerate}
      \def\labelenumiv{3.3.1.2.\arabic{enumiv}.}
      \item
        Nhóm dữ liệu chuẩn học thuật: Fashion-MNIST và CIFAR-10
      \item
        Nhóm dữ liệu ứng dụng thực tiễn: PlantVillage (bệnh lá cây) và Chest-Xray-Pneumonia (viêm phổi)
      \end{enumerate}
    \item
      Các phương pháp Tiền xử lý dữ liệu và Tăng cường dữ liệu (Data Augmentation) cho từng nhóm ảnh
    \item
      Hiện thực hóa 3 mô hình đối chứng thực nghiệm
      \begin{enumerate}
      \def\labelenumiv{3.3.1.4.\arabic{enumiv}.}
      \item
        Mô hình 1 (SE Baseline): TickNet-small nguyên bản sử dụng chú ý SE trong khối
      \item
        Mô hình 2 (CBAM-Local): Chèn trực tiếp CBAM thay SE cục bộ trong khối FR-PDP
      \item
        Mô hình 3 (CBAM-Hook - Đề xuất cải tiến): Chèn CBAM phân cấp tại các điểm nối và cuối mạng
      \end{enumerate}
    \item
      Cấu hình quá trình huấn luyện và tối ưu hóa
      \begin{enumerate}
      \def\labelenumiv{3.3.1.5.\arabic{enumiv}.}
      \item
        Bộ tối ưu SGD và các hàm Loss tương ứng
      \item
        Hàm gọi lại tự động giảm tốc độ học (ReduceLROnPlateau) và dừng sớm (EarlyStopping)
      \end{enumerate}
    \end{enumerate}
  \item
    Kết quả thực nghiệm, phân tích và đánh giá
    \begin{enumerate}
    \def\labelenumiii{3.3.2.\arabic{enumiii}.}
    \item
      Các độ đo hiệu năng chuẩn mực (Accuracy, Precision, Recall, F1-score) và cấu trúc Confusion Matrix
    \item
      Kết quả thực nghiệm và biểu đồ Loss/Accuracy trên tập dữ liệu chuẩn (Fashion-MNIST, CIFAR-10)
    \item
      Kết quả thực nghiệm và biểu đồ Loss/Accuracy trên tập dữ liệu thực tiễn (PlantVillage, Chest-Xray-Pneumonia)
    \item
      Phân tích hiện tượng suy giảm hiệu năng cục bộ (CBAM-Local) và ưu thế của giải pháp đề xuất (CBAM-Hook)
    \item
      Giải thích mô hình bằng trực quan hóa bản đồ nhiệt Grad-CAM (Sự tập trung vùng không gian đặc trưng thực tế)
    \end{enumerate}
  \end{enumerate}
\end{enumerate}

\begin{quote}
CHƯƠNG 4: KẾT QUẢ, KẾT LUẬN VÀ KIẾN NGHỊ
\end{quote}

\begin{enumerate}
\def\labelenumi{4.\arabic{enumi}.}
\item
  KẾT QUẢ ĐẠT ĐƯỢC CỦA ĐỀ TÀI
\item
  KẾT LUẬN (Đánh giá mức độ hoàn thành so với mục tiêu và ý nghĩa thực tiễn)
\item
  KIẾN NGHỊ VÀ CÁC HƯỚNG PHÁT TRIỂN TIẾP THEO
\end{enumerate}

\begin{quote}
TÀI LIỆU THAM KHẢO
\end{quote}

\begin{quote}
\textbf{MỞ ĐẦU}
\end{quote}

\subsection{1. LÝ DO CHỌN ĐỀ TÀI}\label{1-ly-do-chon-de-tai}

\begin{quote}
Phân loại hình ảnh (image classification) là một trong những bài toán cơ
bản trong thị giác máy tính. Bài toán này được phát biểu như sau: cho
một hình ảnh, cần xác định hình ảnh này thuộc lớp nào trong một số lớp
cho trước. Ngày nay nó được áp dụng rộng rãi trong hầu hết các lĩnh vực
như quản lý chất lượng sản phẩm, ô tô tự lái, y tế, \ldots vì vậy việc
nâng cao tính chính xác của việc phân loại hình ảnh là một trong những
vấn đề cấp thiết.

Trong TickNet là một kiến trúc CNN nhẹ được thiết kế nhằm giảm số lượng
tham số và chi phí tính toán nhưng vẫn duy trì hiệu quả phân loại cao
thông qua khối Full-Residual Point-Depth-Point (FR-PDP) và cấu trúc
Tick-shape Backbone. Tuy nhiên, mô hình TickNet gốc chưa khai thác cơ
chế Attention để tăng cường khả năng tập trung vào các đặc trưng quan
trọng của ảnh.

Trong số các cơ chế Attention, Convolutional Block Attention Module
(CBAM) là một mô-đun nhẹ, có khả năng cải thiện việc trích xuất đặc
trưng bằng cách kết hợp chú ý theo kênh và chú ý theo không gian. Vì
vậy, việc nghiên cứu tích hợp CBAM vào TickNet-Small nhằm nâng cao hiệu
quả phân loại ảnh là một hướng nghiên cứu có ý nghĩa cả về mặt học thuật
và ứng dụng thực tiễn.

Xuất phát từ những lý do trên, đề tài "Nghiên cứu cải tiến TickNet-Small
bằng cơ chế Attention CBAM cho bài toán phân loại ảnh" được lựa chọn
thực hiện.
\end{quote}

\subsection{2. MỤC TIÊU VÀ PHẠM VI NGHIÊN CỨU}\label{2-muc-tieu-va-pham-vi-nghien-cuu}

\begin{quote}
Mục tiêu của đề tài hướng đến việc nghiên cứu kiến trúc TickNet-Small và
cơ chế Attention CBAM, từ đó đề xuất mô hình TickNet-Small kết hợp CBAM
nhằm nâng cao hiệu quả của bài toán phân loại ảnh. Đồng thời, tiến hành
thực nghiệm trên nhiều bộ dữ liệu khác nhau để đánh giá khả năng cải
thiện hiệu suất của mô hình so với TickNet-Small gốc.

Phạm vi nghiên cứu:
\end{quote}

\begin{itemize}
\item
  Nghiên cứu kiến trúc TickNet, khối FR-PDP và cấu trúc Tick-shape
  Backbone.
\item
  Nghiên cứu cơ chế Attention, đặc biệt là Convolutional Block Attention
  Module (CBAM).
\item
  Xây dựng mô hình TickNet-Small kết hợp CBAM.
\item
  Huấn luyện và đánh giá mô hình trên các tập dữ liệu MNIST, CIFAR-10,
  PlantVillage và ChestXray.
\item
  Quá trình huấn luyện sử dụng bộ tối ưu hóa SGD kết hợp với cấu hình
  hàm gọi lại động bao gồm giảm tốc độ học ReduceLROnPlateau và dừng
  huấn luyện sớm EarlyStopping.
\item
  Đánh giá hiệu năng dựa trên các độ đo tiêu chuẩn: Accuracy, Precision,
  Recall, F1-score cùng trực quan hóa bằng Confusion Matrix và bản đồ
  nhiệt Grad-CAM để giải thích vùng chú ý của mô hình.
\end{itemize}

\subsection{3. Ý NGHĨA KHOA HỌC VÀ THỰC TIỄN CỦA ĐỀ TÀI}\label{3-y-nghia-khoa-hoc-va-thuc-tien-cua-de-tai}

\begin{quote}
Ý nghĩa khoa học:
\end{quote}

\begin{itemize}
\item
  Giải quyết xung đột kiến trúc trong mạng siêu nhẹ: Đề tài đóng góp một
  lời giải khoa học sâu sắc cho việc tối ưu hóa hiệu năng các mạng
  nơ-ron siêu nhẹ bằng cách chỉ ra và xử lý thành công sự chồng chéo
  thông tin giữa phép tích chập chiều sâu và mô-đun chú ý không gian
  trong các cấu trúc Perceptron cục bộ.
\item
  Đề xuất mô hình chú ý phân cấp mới: Nghiên cứu tiên phong trong việc
  thiết lập cấu trúc phối hợp hai tầng cơ chế chú ý: giữ nguyên
  attention kênh SE ở cấp độ khối cơ bản để lọc đặc trưng cục bộ, và áp
  dụng attention hỗn hợp CBAM ở cấp độ điểm nối để kiểm soát đặc trưng
  toàn cục. Điều này mở ra một hướng tiếp cận mới mẻ trong việc thiết kế
  và tích hợp các khối attention vào các mạng nơ-ron tối giản.
\item
  Góp phần phát triển lý thuyết học sâu giải thích được : Thông qua việc
  ứng dụng thuật toán Grad-CAM để trực quan hóa các bản đồ nhiệt chú ý ,
  đề tài cung cấp các minh chứng khoa học trực quan, rõ ràng giải thích
  tại sao mô hình cải tiến có thể hội tụ nhanh hơn, định vị chính xác
  hơn các vùng bệnh lý hoặc vật thể so với mô hình baseline thông
  thường\ldots{}
\end{itemize}

\begin{quote}
Ý nghĩa thực tiễn: kết quả nghiên cứu trình bày trong tiểu luận có ý
nghĩa rất lớn khi áp dụng vào các bài toán trong thực tiễn như sau
\end{quote}

\begin{itemize}
\item
  Nâng cao chất lượng chẩn đoán y khoa và quản lý nông nghiệp: Kết quả
  thực nghiệm của mô hình cải tiến trên tập dữ liệu Chest-Xray-Pneumonia
  giúp tối ưu hóa độ chính xác và giảm thiểu tỷ lệ bỏ sót trong việc
  chẩn đoán viêm phổi qua hình ảnh X-quang phổi . Đồng thời, thực nghiệm
  trên tập PlantVillage giúp hỗ trợ bà con nông dân nhận diện sớm và
  chính xác các loại bệnh trên lá cây trồng, góp phần thúc đẩy các giải
  pháp nông nghiệp thông minh, công nghệ cao.
\item
  Tính khả thi cao trong triển khai thực tế trên thiết bị biên: Nhờ kế
  thừa ưu điểm siêu nhẹ của mạng TickNets, mô hình TickNet-small cải
  tiến bằng CBAM phân cấp cực kỳ phù hợp để đóng gói và vận hành trực
  tiếp trên các thiết bị di động cá nhân, máy tính bảng cầm tay, thiết
  bị bay không người lái hay các chip nhúng biên IoT có năng lượng và
  năng lực tính toán hạn chế mà không cần phụ thuộc vào máy chủ đám mây.
\item
\end{itemize}

\section{CHƯƠNG 1: TỔNG QUAN ĐỀ TÀI}\label{chux1b0ux1a1ng-1-tux1ed5ng-quan-ux111ux1ec1-tuxe0i}

\subsection{1.1. Bối cảnh lĩnh vực phân loại hình ảnh và sự phát triển của học sâu (Deep Learning)}\label{ch1-1-bg}
\begin{quote}
\emph{[Nội dung phần này đang được cập nhật/viết tiếp...]}
\end{quote}

\subsection{1.2. Thách thức về tài nguyên tính toán và sự cần thiết của mạng nơ-ron siêu nhẹ (Light-weight CNNs)}\label{ch1-2-lightweight}
\begin{quote}
Mặc dù các nghiên cứu tích hợp cơ chế chú ý như CBAM đã chứng minh hiệu
quả vượt trội khi kết hợp với các mạng tích chập truyền thống như VGG16
, việc áp dụng cơ chế này lên các kiến trúc mạng siêu nhẹ vẫn đối mặt
với nhiều thách thức lớn về sự tương thích cấu trúc . Các mạng siêu nhẹ
phổ biến hiện nay như MobileNets hay ShuffleNets thường chỉ sử dụng cơ
chế chú ý một hướng SE để tái cấu trúc kênh đặc trưng nhằm tiết kiệm
tham số . Kiến trúc mạng siêu nhẹ cải tiến gần đây là TickNets -- nổi
bật với khối FR-PDP và cấu trúc xương sống co giãn kênh hình dấu tích
tuy đạt hiệu suất cao nhưng cũng mới chỉ dừng lại ở việc tích hợp mặc
định cơ chế attention SE ở mức cục bộ bên trong khối.
\end{quote}

\subsection{1.3. Các nghiên cứu liên quan về cơ chế chú ý (Attention Mechanisms) trong thị giác máy tính}\label{ch1-3-attention}
\begin{quote}
Cơ chế chú ý (Attention Mechanism) là một trong những cơ chế được sử
dụng phổ biến trong các bài toán xử lý ngôn ngữ tự nhiên (NLP), đặc biệt
là trong các mô hình dịch máy. Một trong những thành tựu nổi bật của nó
là sự phát triển của mô hình Transformer, giới thiệu trong bài báo
``Attention is All You Need'' của Vaswani et al. năm 2017 {[}1{]}. Mô
hình Transformer đã thay đổi cách tiếp cận dịch máy và các bài toán xử
lý ngôn ngữ tự nhiên (NLP) nhờ vào cơ chế Attention mạnh mẽ của nó. Góp
phần dẫn đến sự phát triển của các mô hình như BERT (Bidirectional
Encoder Representations from Transformers) và GPT (Generative
Pre-trained Transformer). Những mô hình này đã thiết lập các chuẩn mực
mới trong nhiều bài toán NLP khác nhau, bao gồm dịch máy, hiểu ngôn ngữ
tự nhiên, và tạo văn bản.
\end{quote}

\subsection{1.4. Phát biểu bài toán, giới hạn nghiên cứu và giải pháp đề xuất cải tiến của tiểu luận}\label{ch1-4-problem}
\begin{quote}
Do đó, bài nghiên cứu này tập trung vào giải pháp cải tiến mạng
TickNet-small bằng cơ chế chú ý phân cấp. Phương pháp này đề xuất giữ
nguyên SE cục bộ trong khối FR-PDP, đồng thời chèn CBAM một cách chiến
lược tại các điểm nối mạng giữa hai backbone co giãn và trước lớp Global
Average Pooling cuối mạng , giúp mô hình phát huy tối đa sức mạnh lọc
nhiễu không gian toàn cục mà không làm bùng nổ tham số tính toán.
\end{quote}

\section{CHƯƠNG 2: CƠ SỞ LÝ THUYẾT}\label{chux1b0ux1a1ng-2}

\subsection{2.1. MẠNG NƠ-RON TÍCH CHẬP SIÊU NHẸ TICKNETS:}\label{2-1-ticknets}

\subsubsection{2.1.1. Khái niệm và đặc điểm của họ mạng tích chập siêu nhẹ TickNets}\label{2-1-1-concept}

\begin{quote}
Neural Network được mô phỏng theo bộ não của con người. Có các nút riêng
lẻ tạo thành các lớp trong mạng, giống như các nơ-ron trong não của
chúng ta kết nối các khu vực khác nhau. Đầu vào cho các nút trong một
lớp sẽ có trọng số được gán cho chúng để thay đổi ảnh hưởng của tham số
đó đến kết quả dự đoán tổng thể. Vì các trọng số được ấn định trên các
liên kết giữa các nút nên mỗi nút có thể bị ảnh hưởng bởi nhiều trọng
số. Mạng nơ-ron lấy tất cả dữ liệu huấn luyện ở lớp đầu vào. Sau đó, nó
chuyển dữ liệu qua các lớp ẩn, chuyển đổi các giá trị dựa trên trọng số
tại mỗi nút. Cuối cùng nó trả về một giá trị ở lớp đầu ra. {[}2{]}
\end{quote}

\paragraph{Kiến trúc mạng nơ-ron:}\label{kiux1ebfn-truxfac-mux1ea1ng-nux1a1-ron}

\begin{quote}
Một Neural Network thường bao gồm 3 tầng cụ thể như sau:
\end{quote}

\begin{itemize}
\item
  \textbf{Input Layer (Tầng đầu vào):} giá trị các node chính là số
  lượng đặc trưng của dữ liệu đầu vào khi đưa vào mô hình. Chúng ta thấy
  giá trị đầu vào là n thuộc tính/đặc trưng.
\item
  \textbf{Hidden Layer (Tầng ẩn):} có số node ẩn thường không được xác
  định, thường do kinh nghiệm của người thiết kế hoặc qua quá trình thử
  nghiệm nhiều lần mà có được. Tuy nhiên thực tế nếu số lượng node ẩn
  quá nhiều thì mạng sẽ học chậm, còn nếu số node quá ít thì mạng sẽ
  không rút trích đủ các thông tin cần thiết trên các đặc trưng. Từ đó
  hiệu quả của mô hình sẽ không được chính xác. Số lượng các lớp ẩn ở
  đây có thể một hoặc nhiều lớp ẩn tùy thuộc vào tính chất cũng như độ
  phức tạp của dữ liệu.
\item
  \textbf{Output Layer (Tầng đầu ra):} giá trị các số node chính là số
  lượng nhãn đầu ra mà chúng ta mong muốn. Ví dụ như trong tập dữ liệu
  của
\end{itemize}

\begin{quote}
chúng ta có tổng cộng 5 nhãn, thì đầu ra của chúng ta tại lớp này chính
là một lớp ẩn với 5 phần tử tương ứng với 5 nhãn.
\end{quote}

\includegraphics[width=4.60017in,height=3.525in]{./media/image3.jpeg}

\begin{quote}
\emph{Hình 2.1. Minh họa về Neural Network}
\end{quote}

\paragraph{Khái niệm mạng nơ-ron tích chập (CNN):}\label{khuxe1i-niux1ec7m-mux1ea1ng-nux1a1-ron-tuxedch-chux1eadp-cnn}

\begin{quote}
Mạng nơ-ron tích chập là một loại mạng nơ-ron đặc biệt có nhiều lớp. CNN
xử lý dữ liệu có sự sắp xếp giống như lưới, sau đó trích xuất các tính
năng quan trọng. CNN được xây dựng để tận dụng các mối quan hệ không
gian (spatial relationships) trong dữ liệu đầu vào, làm cho chúng đặc
biệt hiệu quả trong các tác vụ liên quan đến thị giác máy tính (Computer
Vision.) Một lợi thế rất lớn của việc sử dụng CNN là bạn không cần thực
hiện nhiều bước xử lý trước đối với hình ảnh.

CNN là một trong những mô hình học sâu phổ biến nhất và được sử dụng phổ
biến trong cộng đồng thị giác máy tính (CV). CNN được dùng trong trong
nhiều bài toán như phân loại ảnh, phân tích video, ảnh MRI, hoặc cho bài
các bài toán của lĩnh vực xử lý ngôn ngữ tự nhiên, và hầu hết đều cho ra
kết quả tốt hơn so với các mô hình khác.

Mạng CNN xử lý dữ liệu đầu vào trong không gian ba chiều: chiều rộng,
chiều cao và chiều sâu. Các Nơ-ron trong mạng không liên kết hoàn toàn
với toàn bộ Nơ-ron kế đến mà chỉ liên kết tới một vùng nhỏ. Cuối cùng,
một tầng đầu ra được tối giản thành vec-tơ của giá trị xác suất, đại
diện cho kết quả dự đoán.
\end{quote}

\paragraph{Kiến trúc mạng CNN:}\label{kiux1ebfn-truxfac-mux1ea1ng-cnn}

\begin{quote}
Một mạng CNN bao gồm tập hợp các tầng cơ bản: convolution layer, pooling
layer và fully connected layer. Ở giữa các tầng convolution và pooling
thường có các hàm kích hoạt phi tuyến (nonliner layer) thường là ReLU
hoặc Tanh. Ảnh khi đưa vào mạng sẽ được lan truyền qua tầng convolution
layer, giá trị tính được từ các tầng convolution sẽ đi qua một hàm kích
hoạt, sau đó giá trị này sẽ được lan truyền qua pooling layer, quá trình
này có thể lặp đi lặp lại nhiều lần trước khi đến tầng fully connected
layer. Cuối cùng ảnh sẽ được lan truyền đến tầng fully connected layer
và đi qua hàm kích hoạt Softmax, thường thì cuối cùng sẽ thu được một
vector chứa xác suất phần trăm thuộc về các lớp đối với các bài toán
phân loại.

\includegraphics[width=5.15848in,height=1.1in]{./media/image4.jpeg}
\end{quote}

\emph{Hình 2.2 Ví dụ về kiến trúc của CNN}

\paragraph{Convolution layer:}\label{convolution-layer}

\begin{quote}
Convolution layer là tầng quan trọng nhất và cũng là tầng đầu tiên của
của mô hình CNN. Tầng này có chức năng chính là phát hiện các đặc trưng
có tính không gian hiệu quả. Trong tầng này có 4 đối tượng chính là: ma
trận đầu vào, bộ filters, và receptive field, feature map. Convolution
layer nhận đầu vào là một ma trận 3 chiều và một bộ filters cần phải
học. Bộ filters này sẽ trượt qua từng vị trí trên bức ảnh để tính tích
chập (convolution) giữa bộ filter và phần tương ứng trên bức ảnh. Phần
tương ứng này trên bức ảnh gọi là receptive field, tức là vùng mà một
neuron có thể nhìn thấy để đưa ra quyết định, và mà trận cho ra bởi quá
trình này được gọi là feature map.

Khi đưa ảnh vào mạng, bộ filter sẽ quét qua toàn bộ ảnh cho nên các đặc
trưng cơ bản của ảnh như là góc, cạnh, màu sắc và texture sẽ được mạng
phát hiện ra bất kể nó nằm ở vị trí nào trong ảnh. Do đó tầng
convolution được xem như là một bộ trích chọn đặc trưng (feature
detector) vì nó có chức năng chính là phát hiện đặc trưng cụ thể của bức
ảnh đầu vào. Khi áp dụng phép tính tích chập cho xử lý hình ảnh, người
ta nhận thấy rằng kỹ thuật tích chập này sẽ giúp biến đổi các thông tin
đầu vào thành các yếu tố đặc trưng (nó tương tự như bộ phát hiện nhằm
phát hiện ra các đặc trưng như cạnh, hướng, ...). Cụ thể hơn, tích chập
sẽ trích xuất đặc trưng của ảnh đầu vào qua các vùng ảnh nhỏ. Các vùng
này được gọi là Local Receptive Field (LRF). Tích chập sẽ tính toán trên
các LRF chồng lấp lên nhau. Độ chồng lắp này phụ thuộc vào hệ số trượt S
(stride) của từng kiến trúc mạng cụ thể. Nếu sử dụng với hệ số trượt S =
α, thì tương ứng LRF (bằng kích thước với kernel) sẽ dịch chuyển α đơn
vị pixel sau mỗi lần tích chập. {[}3{]}

\textbf{Padding:} là kỹ thuật thêm các pixel xung quanh biên của đầu vào
(thường là các giá trị 0) trước khi áp dụng phép tích chập. Mục đích của
padding là điều chỉnh kích thước của đầu vào để đạt được các mục đích cụ
thể. (Ví dụ với đầu vào có kích thước 7x7 với bộ filter có kích thước
3x3 và stride bằng 2 thì kích thước để bộ filter trượt với stride bằng 2
là không phù hợp từ đó sử dụng padding để chèn thêm cho ở viền của đầu
vào những phần tử 0 để bộ filter có thể trượt với kích thước phù hợp).

\textbf{Stride:} là số bước di chuyển của bộ filter qua ảnh đầu vào khi
thực hiện phép tích chập. Nó xác định khoảng cách giữa hai vùng lân cận
được xử lý liên tiếp. Stride ảnh hưởng đến kích thước của đầu ra sau mỗi
lớp.

\textbf{Bias:} là một giá trị được cộng vào sau khi thực hiện phép tích
chập. Nó giúp dịch chuyển (shift) giá trị đầu ra, cho phép mô hình học
các đặc trưng tốt hơn và linh hoạt hơn. Mỗi bộ lọc có một bias riêng, và
nó được học trong quá trình huấn luyện cùng với các trọng số của bộ lọc.

\textbf{Shared Weights:} Trong một lớp tích chập, các bộ lọc (filter) sẽ
có cùng một tập trọng số được áp dụng cho toàn bộ đầu vào. Điều này có
nghĩa là cùng một bộ lọc sẽ trượt qua toàn bộ đầu vào để tạo ra các bản
đồ đặc trưng (feature map). Nhờ vậy, số

lượng tham số cần huấn luyện sẽ giảm đáng kể, giúp mô hình dễ dàng tránh
được hiện tượng overfitting. (Ví dụ chúng ta có một hình ảnh đầu vào 5x5
và một bộ lọc 3x3. Nếu không có chia sẻ trọng số, mỗi vị trí trên hình
ảnh cần một bộ trọng số riêng, điều này sẽ tạo ra số lượng tham số rất
lớn. Với trọng số chia sẻ, bộ lọc 3x3 sẽ sử dụng cùng một tập trọng số
để quét qua toàn bộ hình ảnh, chỉ cần huấn luyện 9 trọng số (cho bộ lọc
3x3) và 1 giá trị bias).

\textbf{Shared Bias}: Tương tự như trọng số chia sẻ, các giá trị bias
cũng được chia sẻ cho toàn bộ đầu vào khi bộ lọc di chuyển. Mỗi bộ lọc
sẽ có một giá trị bias duy nhất, áp dụng cho tất cả các vị trí mà bộ lọc
này hoạt động.

Ví dụ có một ảnh đầu vào có kích thước 6x6 với stride bằng 2, padding
bằng 1 và bộ filter có kích thước 3x3 thì đầu ra sẽ là 4x4 như hình bên
dưới:
\end{quote}

\includegraphics[width=5.10009in,height=2.25in]{./media/image5.jpeg}

\begin{quote}
\emph{Hình 2.3 Ví dụ về Output của Convolution}

Theo đó ta có thể tính kích thước đầu ra thông qua công thức sau:

output size = (input size)+2∗padding-(filter size-1)-1 + 1

stride
\end{quote}

\paragraph{Hàm kích hoạt:}\label{huxe0m-kuxedch-houx1ea1t}

\begin{quote}
Hàm kích hoạt là một hàm số nhận vào một giá trị đầu vào và kết quả là
một giá trị có miền giá trị nằm trên một khoảng (hay nửa khoảng) nào đó.
Một số các hàm kích hoạt phổ biến có thể kể đến đó là Sigmod, Tanh,
ReLU. Hàm kích hoạt rất quan trọng bởi vì nó sẽ tăng khả năng dự đoán
của mạng nơ-ron và giúp mô hình

học được các quan hệ phi tuyến phức tạp tiềm ẩn trong dữ liệu. Thông
thường hàm kích hoạt sử dụng ở giữa các tầng convolution và pooling là
hàm Relu. {[}4{]}
\end{quote}

\includegraphics[width=4.59435in,height=2.35625in]{./media/image6.png}

\begin{quote}
\emph{Hình 2.4 Ví dụ về việc sử dụng hàm kích hoạt ReLU}

Hàm ReLU thường được ưa chuộng vì tính toán đơn giản, giúp hạn chế tình
trạng gradient biến mất (vanishing gradient) mà hàm Sigmoid hay tanh
thường gặp, và cũng cho kết quả tốt hơn. Đạo hàm của ReLU là 1 cho các
giá trị lớn hơn 0. Bởi vì nhân 1 với chính nó nhiều lần vẫn cho 1, điều
này về cơ bản giải quyết vấn đề gradient biến mất. Không thể phân biệt
giá trị âm của hàm ReLU vì nó bằng 0. Do đó, các đạo hàm của giá trị âm
chỉ đơn giản được đặt thành 0. Relu cũng như những hàm kích hoạt khác,
được đặt ngay sau tầng convolution.

\includegraphics[width=4.19906in,height=3.0375in]{./media/image7.jpeg}

\emph{Hình 2.5 Đồ thị và công thức của một số hàm kích hoạt}
\end{quote}

\paragraph{Pooling layer:}\label{pooling-layer}

\begin{quote}
Lớp Pooling được sử dụng ngay sau lớp convulational và hàm kích hoạt phi
tuyến. Pooling giúp cho mạng giảm số lượng tham số, từ đó giúp đơn giản
hóa quá trình tính toán của CNN và qua đó góp phần giải quyết vấn đề
overfiting khi huấn luyện mạng và giảm thời gian huấn luyện.

Có nhiều loại pooling nhưng được sử dụng phổ biến nhất là Max-pooling và
Average-pooling. Max-pooling hoạt động bằng cách thực hiện lấy giá trị
lớn nhất trong kích thước Filter mà chúng ta xét. Điều này giúp giữ lại
các đặc trưng mạnh mẽ nhất trong mỗi vùng và loại bỏ các giá trị nhỏ
hơn. Max-pooling giúp làm cho mạng trở nên ít nhạy cảm hơn với các biến
dạng và dịch chuyển nhỏ trong ảnh đầu vào, bởi vì các đặc trưng mạnh mẽ
vẫn được giữ lại dù có sự thay đổi nhỏ. Nhược điểm của Max-pooling là có
thể làm mất đi một số thông tin quan trọng. Còn Average-pooling hoạt
động bằng bằng cách lấy trung bình cộng của tất cả các giá trị của
Filter. Giúp giữ lại nhiều thông tin hơn từ đặc trưng gốc. Có thể giúp
mô hình học được các mối quan hệ phức tạp giữa các đặc điểm. Nhược điểm
của Average-pooling là có thể khiến mô hình trở nên tốn kém về mặt tính
toán hơn so với Max-pooling.

\includegraphics[width=5.03102in,height=4.21667in]{./media/image8.jpeg}

\emph{Hình 2.4 Ví dụ về Max-pooling và Average-pooling}
\end{quote}

\paragraph{Lớp Normalization:}\label{lux1edbp-normalization}

\begin{quote}
Lớp Normalization (Norm) là lớp giúp chuẩn hóa dữ liệu đầu ra cho các
lớp trong CNN trước khi được truyền đi tiếp. Trong những kiến trúc CNN
lớn và phức tạp, lớp Norm sẽ chuẩn hóa các giá trị neural trước khi
chúng được truyền đến hàm Relu. Hàm Relu tuy giúp rút ngắn thời gian
huấn luyện, nhưng nếu không điều chỉnh trọng số phù hợp, hàm Relu sẽ rất
dễ gặp phải vấn đề "dying Relu" khiến cho mạng trở nên chậm hơn khi huấn
luyện. Lớp Norm lúc này sẽ chuẩn hóa và tạo ra các giá trị tích chập phù
hợp để tránh cho Relu rơi vào giá trị 0. Tránh việc gradient xấp xỉ bằng
0 khiến cho tốc độ học của mạng trở nên rất chậm. Các Norm thường dùng
bao gồm Batch Normalization và Layer Normalization. Batch Normalization
là lựa chọn phổ biến cho các mô hình CNN với batch size lớn. Layer
Normalization thích hợp cho các mô hình tuần tự và khi batch size nhỏ.
\end{quote}

\paragraph{Fully connected Layer:}\label{fully-connected-layer}



\begin{quote}
Tầng cuối cùng của mô hình CNN trong bài toán phân loại ảnh là tầng
fully connected layer. Sau khi ảnh được truyền qua nhiều convolutional
layer và pooling layer thì model đã học được tương đối các đặc điểm của
ảnh thì tensor của output của layer cuối cùng sẽ được là phẳng thành
vector và đưa vào một lớp được kết nối như một mạng nơ-ron. Với FC layer
được kết hợp với các tính năng lại với nhau để tạo ra một mô hình. Cuối
cùng sử dụng softmax hoặc sigmoid để phân loại đầu ra . Đầu ra là một
vec-tơ chứa xác xuất của các đối tượng cần dự đoán. Ví dụ trong một bài
toán phân loại ảnh với có 10 lớp, tầng fully connected layer sẽ chuyển
ma trận đặc trưng của tầng trước thành vector có 10 chiều thể hiện xác
suất của 10 lớp tương ứng.
\end{quote}

\pandocbounded{\includegraphics[keepaspectratio]{./media/image10.jpeg}}

\begin{quote}
\emph{Hình 2.5 Minh họa Fully Connected Layer}
\end{quote}

\paragraph{Mô hình VGG16:}\label{muxf4-huxecnh-vgg16}

\begin{quote}
Mô hình VGG16 là một kiến trúc của mạng CNN nổi tiếng, được phát triển
bởi nhóm nghiên cứu tại Visual Geometry Group (VGG) thuộc Đại học Oxford
{[}5{]}. VGG16 thường được sử dụng cho các nhiệm vụ phân loại hình ảnh.
Mạng bao gồm 16 lớp nơ-ron, mỗi lớp hoạt động để xử lý thông tin hình
ảnh tăng dần và cải thiện độ chính xác của dự đoán. Mặc dù đơn giản so
với các kiến trúc gần đây hơn, VGG16 vẫn là lựa chọn phổ biến cho nhiều
ứng dụng deep learning do tính linh hoạt và hiệu suất tuyệt vời của nó.
Mạng VGG16 được đào tạo trên hơn một triệu hình ảnh từ cơ sở dữ liệu
trực quan ImageNet và có thể phân loại hình ảnh thành 1.000 lớp khác
nhau với độ chính xác trong thử nghiệm top 5 là 92,7\%. VGG16 có thể
được áp dụng để xác định xem hình ảnh có chứa một số vật phẩm, động vật,
thực vật và hơn thế nữa.

VGG16 có tổng cộng 16 lớp chính (do đó có tên là VGG16, bao gồm 13 lớp
tích chập (convolutional layers) và 3 lớp kết nối đầy đủ (fully
connected layers). Các khối này được tổ chức thành các khối, mỗi khối
bao gồm các lớp tích chập (convolutional layer) và 1 lớp pooling. Mỗi
khối tích chập sử dụng các bộ filter kích thước 3x3 với số lượng kênh
đầu ra tăng gấp đôi sau mỗi khối. Bắt đầu với 64 bộ filter cho khối đầu
tiên và tăng lên 128, 256, 512, sau đó là 512 cho các khối tiếp theo. Sử
dụng hàm kích hoạt ReLU sau mỗi lớp tích chập. Áp dụng lớp max pooling
2x2 sau mỗi khối để giảm kích thước không gian của dữ liệu và giảm bớt
thông số mô hình. Sau các khối tích chập, VGG16 sử dụng 3 lớp fully
connected với kích thước lần lượt là 4096, 4096 và 1000. Mỗi lớp kết nối
đầy đủ được theo sau bởi một hàm kích hoạt ReLU và một lớp dropout với
tỷ lệ 0.5 để tránh quá phù hợp. Hàm kích hoạt softmax được áp dụng cho
đầu ra của lớp fully connected thứ ba để phân loại. Lớp đầu ra có 1000
nút, tương ứng với 1000 lớp trong tập dữ liệu ImageNet.
\end{quote}

\includegraphics[width=5.24056in,height=1.5125in]{./media/image11.jpeg}

\begin{quote}
\emph{Hình 2.6 Kiến trúc mô hình VGG16}
\end{quote}

\paragraph{Cách thức hoạt động của mạng CNN:}\label{cuxe1ch-thux1ee9c-houx1ea1t-ux111ux1ec9ng-cux1ee7a-mux1ea1ng-cnn}

\begin{quote}
Ví dụ với đầu vào là một ảnh có kích thước 3x3 với 3 kênh màu và bộ
filter có kích thước 2x2 với kênh màu tương ứng như hình sau:

\includegraphics[width=5.39764in,height=4.125in]{./media/image12.jpeg}

\emph{Hình 2.7 Ví dụ cách thức hoạt động của CNN}

Với đầu vào là một bức ảnh có kích thước là 3x3 với 3 kênh màu được biểu
diễn dưới dạng ma trận có kích thước tương ứng, với mỗi phần tử trong ma
trận tương ứng với giá trị cường độ pixel của ảnh tại vị trí đó và bộ
filter chứa trọng số cũng được biểu diễn dưới dạng ma trận 2x2 có cùng
kênh màu. Với Stride = 1 và Padding = 0, đầu tiên bộ filter sẽ trượt
trên bức ảnh và thực hiện tích chập bằng cách di chuyển qua từng vị trí
trên ảnh đầu vào và thực hiện phép nhân điểm tương ứng giữa các phần tử
của bộ lọc và vùng ảnh tương ứng, sau đó cộng lại để tạo ra một giá trị
mới trong bản đồ đặc trưng, sau đó tiếp tục di chuyển bộ filter đến vị
trí tiếp theo dựa trên Stride, thực hiện tiếp tục phép tích chập đi qua
toàn bộ ảnh, sau khi thực hiện phép tích chập, ta thêm vào một giá trị
bias. Sau đó cho đi qua một hàm kích hoạt để thu được một future map
chứa các kết quả đã tính toán. Ta có một công thức tính toán cho output
của phép tích chập như sau:

\includegraphics[width=1.49956in,height=0.3in]{./media/image13.jpeg}

\emph{Hình 2.8 Công thức tính toán output của Convolution}

Trong đó:
\end{quote}

\begin{itemize}
\item
  y là output của phép tích chập.
\item
  x là input.
\item
  f là filter.
\item
  b là bias
\end{itemize}

\begin{quote}
Ví dụ theo Hình 11. ta có thể tính được giá trị của y11 như sau:

y11 = 𝑥\textsuperscript{1} * 𝑘\textsuperscript{1} + 𝑥\textsuperscript{1}
* 𝑘\textsuperscript{1} + 𝑥\textsuperscript{1} * 𝑘\textsuperscript{1} +
𝑥\textsuperscript{1} * 𝑘\textsuperscript{1}
\end{quote}

11 11

\begin{quote}
12 12 21 21

22 22

+ 𝑥\textsuperscript{2} * 𝑘\textsuperscript{2} + 𝑥\textsuperscript{2} *
𝑘\textsuperscript{2} + 𝑥\textsuperscript{2} * 𝑘\textsuperscript{2} +
𝑥\textsuperscript{2} * 𝑘\textsuperscript{2}
\end{quote}

11 11

\begin{quote}
12 12 21 21

22 22

+ 𝑥\textsuperscript{3} * 𝑘\textsuperscript{3} + 𝑥\textsuperscript{3} *
𝑘\textsuperscript{3} + 𝑥\textsuperscript{3} * 𝑘\textsuperscript{3} +
𝑥\textsuperscript{3} * 𝑘\textsuperscript{3}

11 11 12 12 21 21 22 22

+ 𝑏11

Và các y12, y21, y22 được tính toán theo cách tương tự.

Sau khi thực hiện tính toán ở tầng tích chập, dữ liệu sẽ được chuyển
tiếp đến tầng pooling để thực hiện giảm kích thước của các future map.
Thông thường sử dụng Max-pooling hoặc Average-pooling để giảm bớt thông
số và tăng tính tương tác của mạng. Với Max-pooling thì lấy giá trị lớn
nhất trong vùng đang xét của bộ filter, còn Average-pooling thì lấy giá
trị trung bình. Lớp pool sử dụng nhiều bộ filter khác nhau để xác định
các phần khác nhau của hình ảnh như cạnh, góc, thân, lông, mắt và mỏ của
ảnh. Quá trình Convolution + ReLU -- Pooling có thể lặp đi lặp lại nhiều
lần để tạo ra các future map ở cấp độ cao hơn. Mỗi lần lặp, kích thước
của future map sẽ giảm dần, nhưng số lượng future map có thể tăng lên.
Sau đó các future map sẽ được sẽ được chuyển đổi từ một ma trận nhiều
chiều thành một ma trận một chiều gọi là Flattening và đưa đến tầng
Fully Connected. Tiếp theo từ dữ liệu làm phẳng, mà mỗi đầu vào đó được
kết nối đến tất cả nơ-ron để tính toán các giá trị đầu ra dựa trên các
trọng số và bias đã học. Thông thường, khi tất cả các tính năng được kết
nối với lớp

FC, nó có thể gây ra hiện tượng overfitting trong tập dữ liệu huấn
luyện. Overfitting xảy ra khi một mô hình cụ thể hoạt động quá tốt trên
dữ liệu huấn luyện gây ra tác động tiêu cực đến hiệu suất của mô hình
khi được sử dụng trên dữ liệu mới. Để giải quyết vấn đề này một lớp
``Dropout'' được thêm vào, trong đó một số nơ-ron bị loại khỏi mạng
trong quá trình huấn luyện để giảm kích thước mô hình. Sau khi thông qua
các lớp FC, đầu ra của lớp cuối cùng có nơ-ron sẽ bằng với số tổng lớp
cần phân đối với bài toán phân loại hình ảnh, và đi qua một hàm kích
hoạt Softmax hoặc Sigmoid để phân bổ xác xuất cuối cùng. Lớp có xác suất
cao nhất là lớp cuối cùng của ảnh đầu vào.
\end{quote}

\paragraph{Cách chọn tham số cho CNN: {[}6{]}}\label{cuxe1ch-chux1ecdn-tham-sux1ed1-cho-cnn-6}

\begin{itemize}
\item
  Số các convolution layer: càng nhiều các convolution layer thì hiệu
  suất càng được cải thiện. Sau khoảng 3 hoặc 4 layer, các tác động được
  giảm một cách đáng kể
\item
  Filter size: thường filter theo size 5×5 hoặc 3×3.
\item
  Pooling size: thường là 2×2 hoặc 4×4 cho ảnh đầu vào lớn.
\end{itemize}

\subsubsection{2.1.2. Các phép tích chập tối giản (Depthwise Convolution và Pointwise Convolution)}\label{2-1-2-dwconv}
\begin{quote}
\emph{[Nội dung phần này đang được cập nhật/viết tiếp...]}
\end{quote}

\subsubsection{2.1.3. Kiến trúc khối Full-Residual Point-Depth-Point (FR-PDP)}\label{2-1-3-frpdp}
\begin{quote}
\emph{[Nội dung phần này đang được cập nhật/viết tiếp...]}
\end{quote}

\subsubsection{2.1.4. Cấu trúc xương sống co giãn kênh hình dấu tích (Tick-shape Backbone) và mô hình TickNet-small}\label{2-1-4-backbone}
\begin{quote}
\emph{[Nội dung phần này đang được cập nhật/viết tiếp...]}
\end{quote}

\subsubsection{2.1.5. Cơ chế lan truyền đặc trưng và cách thức hoạt động của TickNets}\label{2-1-5-propagation}
\begin{quote}
\emph{[Nội dung phần này đang được cập nhật/viết tiếp...]}
\end{quote}

\subsection{2.2. CƠ CHẾ CHÚ Ý (ATTENTION MECHANISM):}\label{cux1a1-chux1ebf-chuxfa-uxfd-attention-mechanism}

\begin{quote}
Cơ chế chú ý trong Neural Network có xu hướng bắt chước sự chú ý về mặt
nhận thức của con người. Một đặc tính quan trọng của thị giác con người
là người ta không cố gắng xử lý toàn bộ cảnh vật trong cùng một lúc.
Thay vào đó, con người khai thác chuỗi những cái nhìn thoáng qua một
phần và tập trung có chọn lọc vào những điểm nổi bật các bộ phận để nắm
bắt cấu trúc hình ảnh tốt hơn. Đó cũng là ý tưởng chính cho cơ chế này
nhằm nhấn mạnh những phần quan trọng của thông tin và cố gắng giảm bớt
những phần không liên quan. Vì bộ nhớ làm việc ở cả con người và máy móc
đều có hạn nên quá trình này là chìa khóa để không làm quá tải bộ nhớ
của hệ thống. Trong học sâu, sự chú ý có thể được hiểu là một vectơ có
trọng số quan trọng.

Trong thị giác máy tính (CV), cơ chế chú ý cho phép mô hình tập trung
vào các vùng cụ thể của hình ảnh và loại bỏ đi những vùng không quan
trọng. Ở mức đơn giản, cơ chế Attention trong phân loại hình ảnh hoạt
động bằng cách tạo ra một trọng

số cho từng phần của hình ảnh đầu vào. Các trọng số này cho biết mức độ
quan trọng của mỗi phần vùng hình ảnh đối với quá trình phân loại. Trong
quá trình huấn luyện, mô hình học từ dữ liệu để xác định các vùng quan
trọng này. {[}7{]}
\end{quote}

\subsubsection{2.2.1. Cơ chế chú ý theo kênh Squeeze-and-Excitation (SE) mặc định trong TickNets}\label{2-2-1-se-attention}
\begin{quote}
\emph{[Nội dung phần này đang được cập nhật/viết tiếp...]}
\end{quote}

\subsubsection{2.2.2. Mô-đun chú ý theo kênh (Channel Attention Module - CAM) của CBAM}\label{channel-attention-module-cam}

\begin{quote}
Channel Attention Module (CAM) là một thành phần trong các mô hình mạng
nơ-ron sâu, đặc biệt là trong lĩnh vực xử lý ảnh và thị giác máy tính,
nhằm cải thiện khả năng chú ý của mạng nơ-ron tới các kênh thông tin
quan trọng trong một tập hợp các đặc trưng. CAM được thiết kế để gán
trọng số khác nhau cho các kênh đặc trưng (feature channel), giúp mạng
nơ-ron có thể tập trung vào các kênh quan trọng hơn trong quá trình xử
lý dữ liệu. Điều này cải thiện khả năng của mạng nơ-ron trong việc phát
hiện và phân loại các đặc trưng có ý nghĩa hơn từ dữ liệu đầu vào.

CAM hoạt động bằng cách tính toán một bản đồ trọng số chú ý (attention
map) cho từng kênh đặc trưng. Mỗi kênh của feature map được coi như là
một feature detector, channel attention tập trung vào ``điều gì'' có ý
nghĩa đối với ảnh đầu vào. Quá trình hoạt động của CAM thường bao gồm
các bước sau theo hình sau:
\end{quote}

\includegraphics[width=5.20419in,height=1.72656in]{./media/image14.png}

\begin{quote}
\emph{Hình 2.9 Mô hình hoạt động của CAM}

CAM nhận đầu vào là một input feature 𝐹 có kích thước 𝐻×𝑊×𝐶, với 𝐻 là
chiều cao, W là chiều rộng và C là số kênh. Sau đó thực hiện thu gom
thông tin thông tin toàn cục bằng cách sử dụng kết hợp cả hoạt động
Global Max-pooling và Global

Average-pooling để tạo ra hai bộ vector đặc trưng khác nhau là: Favg và
Fmax. Sau đó hai bộ vector đại diện được đưa qua một mạng con nhỏ
multi-layer perceptron (MLP). Thường thì mạng con này gồm hai lớp fully
connected với một hàm kích hoạt phi

tuyến ReLU ở giữa. Để giảm số lượng tham số cần học, kích thước của lớp
kích hoạt ẩn được đặt bằng R\textsuperscript{C/r×1×1}, với r là tỉ lệ
giảm. Kết quả của MLP được kết hợp bằng cách sử dụng phép tính tổng theo
phần tử và chuyển qua hàm kích hoạt sigmoid để tạo ra một vector trọng
số gọi là channel attention map MC(F) ∈ R \textsuperscript{Cx1x1}. Công
thức tính toán cho quá trình đó được biểu diễn như sau:

MC(F) = σ (MLP (Avgpool(F)) + MLP(Maxpool(F)))

= σ (W1 (W0 (Favg )) + W1 ( W0 (Fmax)))

Trong đó :

σ : là hàm kích hoạt sigmod.

W1, W0: lần lượt là trọng số của MLP.
\end{quote}

\subsubsection{2.2.3. Mô-đun chú ý theo không gian (Spatial Attention Module - SAM) của CBAM}\label{spatial-attention-module-sam}

\begin{quote}
Spatial Attention Module (SAM) được sử dụng để cải thiện hiệu suất của
mô hình bằng cách tập trung vào các vị trí quan trọng trong feature map.
Khác với Channel Attention Module (CAM) tập trung vào các kênh, SAM tập
trung vào không gian của các feature maps, xác định các vùng quan trọng
trong ảnh dựa vào quan hệ của chúng, nó cho phép mô hình học cách chú ý
đến các vùng không gian cụ thể trong dữ liệu có liên quan nhất đến nhiệm
vụ đang thực hiện.

SAM hoạt động bằng cách tạo ra một bản đồ chú ý không gian (spatial
attention map), giúp xác định và nhấn mạnh các vùng quan trọng trong ảnh
đầu vào, dựa trên các đặc điểm không gian. Điều này giúp mô hình tập
trung nhiều hơn vào các vùng quan trọng và bỏ qua những vùng không quan
trọng, từ đó cải thiện khả năng nhận dạng và phân loại. Quá trình SAM
thực hiện theo hình sau:

\includegraphics[width=5.66201in,height=2.86667in]{./media/image15.jpeg}

\emph{Hình 2.10 Mô hình hoạt động của SAM}

SAM thường bắt đầu bằng việc tổng hợp thông tin từ tất cả các kênh trong
đầu vào để tạo ra một bản đồ không gian duy nhất. Để tạo ra một bản đồ
không gian duy nhất, thông thường, hai phương pháp phổ biến là Global
Average-pooling (GAP) và Global Max-pooling (GMP) được sử dụng. Các bản
đồ kết quả từ GAP và GMP được kết hợp lại để tạo ra một bản đồ chú ý
không gian duy nhất. Thông thường, các giá trị này được kết hợp bằng
cách nối chúng hoặc thực hiện phép cộng element-wise. Sau đó, bản đồ kết
hợp được đưa qua một lớp convolution với bộ filter có kích thước k×k
(thông thường là 7x7), sau đó đi qua một hàm kích hoạt sigmod, kết quả
là một bản đồ chú ý không gian 2D (spatial attention map) MS(F). Công
thức tính toán bản đồ chú ý không gian là:

MS(F) = σ (f\textsuperscript{7x7} ({[}AvgPool(F);MaxPool(F){]})) Trong
đó:

f\textsuperscript{7x7} : Biểu thị phép toán tích chập với kích thước bộ
lọc 7x7.

σ : Biểu thị hàm kích hoạt sigmod.
\end{quote}

\subsubsection{2.2.4. Mô-đun chú ý khối tích chập CBAM hoạt động tuần tự (CAM + SAM)}\label{convolutional-block-attention-module-cbam}

\begin{quote}
Convolutional Block Attention Module (CBAM) là một kỹ thuật tiên tiến
trong lĩnh vực học sâu, được giới thiệu nhằm cải thiện hiệu suất của các
mạng nơ-ron tích chập (Convolutional Neural Networks - CNN). CBAM tích
hợp các cơ chế chú ý (attention mechanisms) vào các khối tích chập để
giúp mạng nơ-ron tập trung vào các đặc trưng quan trọng của dữ liệu đầu
vào. CBAM là sự kết hợp giữa hai mô-đun là mô-đun chú ý kênh (Channel
Attention Module -- CAM) và mô đun chú ý không gian (Spatial Attention
Module -- SAM). Những mô-đun này hoạt động tuần tự để tăng cường khả
năng trích xuất đặc trưng của mạng nơ-ron.
\end{quote}

\includegraphics[width=5.72685in,height=3.38125in]{./media/image16.jpeg}

\begin{quote}
\emph{Hình 2.11 Mô hình CBAM}

Đầu tiên với input feature F đầu vào sẽ qua CAM thực hiện các phép tính
liên quan cho ra bản đồ chú ý không gian (channel attention map) Mc(F).
Sau đó Mc(F) được nhân theo từng phần tử với input future ban đầu F thu
được future map F'. Sau đó future map F' có được sau khi đi qua CAM sẽ
được đưa vào đầu vào của SAM thực hiện, thu được kết quả là một bản đồ
chú ý không gian 2D (spatial attention map) Ms(F'). Bản đồ chú ý không
gian Ms(F') được nhân theo từng phần tử đầu vào F′ để

cho ra đầu ra cuối cùng F'\,' hay CBAM feature. Quá trình thực hiện tính
toán được biểu diễn dưới dạng công thức như sau:

F' = Mc(F) ⊗F F'\,' = Ms(F') ⊗ F'

Trong đó:

⊗ : là phép nhân theo từng phần tử (element-wise multiplication).

Một trong những lợi ích chính của CBAM đặc biệt là khả năng áp dụng bản
đồ chú ý dọc theo cả kênh và chiều không gian. Điều này cho phép mô hình
tập trung hiệu quả hơn vào cả các tính năng cấp cao (chẳng hạn như đối
tượng hoặc hình dạng) và các tính năng cấp thấp (chẳng hạn như các cạnh
hoặc kết cấu) cùng một lúc. Điều này có thể đặc biệt hữu ích trong các
tác vụ như phát hiện đối tượng, trong đó cả hai loại tính năng đều quan
trọng để xác định chính xác đối tượng. {[}9{]}{[}10{]}{[}11{]}
\end{quote}

\subsection{2.3. PHƯƠNG PHÁP NGHIÊN CỨU}\label{phuong-phap-nghien-cuu}

Đề tài được triển khai dựa trên sự kết hợp giữa phương pháp nghiên cứu lý thuyết và phương pháp nghiên cứu thực nghiệm.

\textbf{Về phương pháp lý thuyết,} nghiên cứu tiến hành khảo cứu hệ thống các công trình khoa học liên quan bao gồm bài báo gốc về kiến trúc TickNets, cơ chế chú ý CBAM (Woo et al., 2018) và cơ chế Squeeze-and-Excitation (Hu et al., 2018). Trên cơ sở đó, đề tài phân tích cấu trúc luồng đặc trưng trong khối FR-PDP, vai trò của các phép tích chập tối giản Depthwise và Pointwise, cũng như cơ chế hoạt động tuần tự của CBAM. Từ kết quả phân tích lý thuyết, nghiên cứu xây dựng các giả thuyết khoa học về mối quan hệ giữa vị trí chèn cơ chế chú ý và hiệu năng mô hình, đặc biệt là hiện tượng chồng chéo đặc trưng không gian khi đặt Spatial Attention Module cùng mức cục bộ với Depthwise Convolution bên trong khối FR-PDP.

\textbf{Về phương pháp thực nghiệm,} nghiên cứu áp dụng phương pháp thiết kế thí nghiệm đối chứng với 3 mô hình (SE Baseline, CBAM-Local và CBAM-Hook) được huấn luyện trên cùng điều kiện để kiểm chứng các giả thuyết đã đặt ra. Bốn bộ dữ liệu được lựa chọn đại diện cho hai nhóm bài toán: nhóm chuẩn học thuật (Fashion-MNIST, CIFAR-10) và nhóm ứng dụng thực tiễn (PlantVillage, Chest-Xray-Pneumonia). Hiệu năng của các mô hình được đánh giá thông qua bộ độ đo chuẩn mực gồm Accuracy, Precision, Recall, F1-Score và ma trận nhầm lẫn (Confusion Matrix), kết hợp với trực quan hóa bản đồ nhiệt Grad-CAM để giải thích vùng chú ý của mô hình. Toàn bộ quá trình thực nghiệm được triển khai trên nền tảng Kaggle Notebooks sử dụng ngôn ngữ Python với framework TensorFlow/Keras. Chi tiết về môi trường, dữ liệu, cấu hình huấn luyện và kết quả thực nghiệm được trình bày đầy đủ tại Chương 3.

\section{CHƯƠNG 3: BÀI TOÁN PHÂN LOẠI HÌNH ẢNH VÀ ĐỀ XUẤT CẢI TIẾN PHÂN CẤP}\label{chux1b0ux1a1ng-3}

\subsection{3.1. GIỚI THIỆU BÀI TOÁN PHÂN LOẠI HÌNH ẢNH}\label{gioi-thieu-bai-toan}

\subsubsection{3.1.1. Phát biểu bài toán}

\begin{quote}
Phân loại hình ảnh (Image Classification) là một trong những bài toán nền tảng và quan trọng nhất trong lĩnh vực thị giác máy tính (Computer Vision). Bài toán được phát biểu một cách hình thức như sau: cho một hình ảnh đầu vào $I$ và một tập hữu hạn các nhãn lớp $\mathcal{C} = \{c_1, c_2, \ldots, c_K\}$ đã được xác định trước, hệ thống cần xây dựng một hàm ánh xạ $f: I \rightarrow \mathcal{C}$ để gán nhãn lớp phù hợp nhất cho hình ảnh đó. Nói cách khác, mục tiêu là xác định hình ảnh thuộc vào \textbf{nhóm nào} trong số các nhóm đã được định nghĩa sẵn.

Ví dụ minh hoạ cụ thể: trong một hệ thống nhận dạng động vật, tập nhãn có thể bao gồm $\mathcal{C} = \{\text{chó, mèo, chim, ngựa}\}$. Khi nhận một bức ảnh chụp một con chó Golden Retriever, hệ thống cần phân tích các đặc trưng thị giác của ảnh --- bao gồm hình dạng, kết cấu lông, tỉ lệ cơ thể --- để đưa ra dự đoán chính xác rằng ảnh thuộc lớp ``chó'' với xác suất tin cậy cao nhất.
\end{quote}

\begin{figure}[htbp]
\centering
\includegraphics[width=0.85\textwidth]{./media/image_classification_overview.png}
\caption{Tổng quan bài toán phân loại ảnh: từ ảnh đầu vào qua mô hình CNN đến kết quả dự đoán xác suất các lớp}
\label{fig:classification-overview}
\end{figure}

\subsubsection{3.1.2. Biểu diễn hình ảnh trong máy tính}

\begin{quote}
Đối với con người, việc nhận dạng nội dung một bức ảnh diễn ra gần như tức thì và tự nhiên. Tuy nhiên, đối với máy tính, một hình ảnh chỉ đơn thuần là một mảng ba chiều (tensor) chứa các giá trị số nguyên. Cụ thể, một ảnh màu RGB được biểu diễn dưới dạng tensor $\mathbf{X} \in \mathbb{R}^{H \times W \times C}$, trong đó $H$ là chiều cao, $W$ là chiều rộng (tính bằng số pixel), và $C = 3$ tương ứng với ba kênh màu Đỏ (Red), Xanh lá (Green), Xanh dương (Blue). Mỗi phần tử trong tensor nhận giá trị từ 0 đến 255, thể hiện cường độ sáng của pixel tại vị trí tương ứng trên kênh màu đó.
\end{quote}

\begin{figure}[htbp]
\centering
\includegraphics[width=0.80\textwidth]{./media/image_representation.png}
\caption{So sánh cách con người nhìn ảnh (trực quan) và cách máy tính biểu diễn ảnh (tensor $H \times W \times C$ các giá trị pixel)}
\label{fig:image-representation}
\end{figure}

\begin{quote}
Ví dụ, một bức ảnh kích thước $224 \times 224$ pixel với 3 kênh màu sẽ chứa tổng cộng $224 \times 224 \times 3 = 150.528$ giá trị số. Nhiệm vụ của thuật toán phân loại chính là từ tập hợp hàng trăm nghìn con số này, rút trích ra được các \textbf{đặc trưng có ý nghĩa ngữ nghĩa} (semantic features) để phân biệt giữa các lớp đối tượng khác nhau. Khoảng cách lớn giữa thông tin mức pixel thô và ý nghĩa ngữ nghĩa cấp cao này được gọi là \textbf{khoảng trống ngữ nghĩa} (semantic gap), và đây chính là thách thức cốt lõi mà các phương pháp phân loại ảnh cần giải quyết.
\end{quote}

\subsubsection{3.1.3. Các thách thức chính của bài toán}

\begin{quote}
Mặc dù bài toán phân loại hình ảnh có thể được phát biểu một cách đơn giản, việc xây dựng một hệ thống phân loại đạt độ chính xác cao trong thực tế đối mặt với nhiều thách thức nghiêm trọng:
\end{quote}

\begin{figure}[htbp]
\centering
\includegraphics[width=0.80\textwidth]{./media/classification_challenges.png}
\caption{Bốn thách thức chính của bài toán phân loại ảnh: biến đổi góc nhìn, biến đổi tỉ lệ, che khuất và đa dạng nội lớp}
\label{fig:classification-challenges}
\end{figure}

\begin{itemize}
\item \textbf{Biến đổi góc nhìn (Viewpoint Variation):} Cùng một đối tượng khi được chụp từ các góc độ khác nhau sẽ tạo ra các hình ảnh có hình dạng rất khác biệt trên mặt phẳng 2D. Ví dụ, một chiếc ô tô chụp từ phía trước trông hoàn toàn khác so với khi chụp từ phía trên.

\item \textbf{Biến đổi tỉ lệ (Scale Variation):} Đối tượng có thể xuất hiện ở nhiều kích thước khác nhau trong ảnh, từ rất nhỏ (chiếm vài pixel) đến rất lớn (phủ gần toàn bộ khung hình). Hệ thống phải nhận dạng được đối tượng bất kể tỉ lệ.

\item \textbf{Hiện tượng che khuất (Occlusion):} Trong nhiều tình huống thực tế, đối tượng cần nhận dạng bị che khuất một phần hoặc phần lớn bởi các vật thể khác. Khi chỉ một phần nhỏ của đối tượng hiển thị, việc phân loại trở nên rất khó khăn.

\item \textbf{Đa dạng nội lớp (Intra-class Variation):} Các đối tượng thuộc cùng một lớp có thể có hình dạng, màu sắc và kết cấu rất khác nhau. Chẳng hạn, lớp ``ghế'' bao gồm ghế sofa, ghế gỗ, ghế xoay văn phòng --- tất cả đều là ghế nhưng có hình dáng bên ngoài rất khác biệt.

\item \textbf{Biến đổi điều kiện chiếu sáng (Illumination Variation):} Cùng một đối tượng dưới các điều kiện ánh sáng khác nhau (ngoài trời, trong nhà, ban đêm) sẽ tạo ra các giá trị pixel rất khác nhau, gây khó khăn cho việc nhận dạng.

\item \textbf{Nhiễu nền phức tạp (Background Clutter):} Khi đối tượng có màu sắc hoặc kết cấu tương tự với nền xung quanh, ranh giới giữa đối tượng và nền trở nên mờ nhạt, khiến việc tách biệt và phân loại gặp nhiều trở ngại.
\end{itemize}

\subsubsection{3.1.4. Ứng dụng thực tiễn}

\begin{quote}
Phân loại hình ảnh ngày nay đóng vai trò then chốt trong rất nhiều lĩnh vực ứng dụng thực tiễn, mang lại giá trị to lớn cho cả đời sống và sản xuất:
\end{quote}

\begin{figure}[htbp]
\centering
\includegraphics[width=0.80\textwidth]{./media/classification_applications.png}
\caption{Các lĩnh vực ứng dụng chính của phân loại hình ảnh trong thực tiễn}
\label{fig:classification-applications}
\end{figure}

\begin{itemize}
\item \textbf{Y tế và chẩn đoán hình ảnh:} Phân loại ảnh X-quang để phát hiện viêm phổi, ung thư phổi; phân tích ảnh mô bệnh học để hỗ trợ bác sĩ chẩn đoán chính xác hơn; sàng lọc bệnh lý võng mạc từ ảnh quang học đáy mắt. Đặc biệt, bộ dữ liệu \textbf{Chest-Xray-Pneumonia} được sử dụng trong nghiên cứu này chính là một ví dụ điển hình cho ứng dụng phân loại ảnh y tế.

\item \textbf{Nông nghiệp thông minh:} Nhận dạng và phân loại bệnh trên lá cây trồng từ ảnh chụp, giúp nông dân phát hiện sớm dịch bệnh và có biện pháp xử lý kịp thời. Bộ dữ liệu \textbf{PlantVillage} trong tiểu luận này là minh chứng cho tiềm năng ứng dụng của công nghệ phân loại ảnh trong nông nghiệp chính xác.

\item \textbf{Xe tự hành:} Nhận dạng biển báo giao thông, phân loại đối tượng trên đường (người đi bộ, xe cộ, chướng ngại vật) là thành phần không thể thiếu trong hệ thống lái xe tự động.

\item \textbf{Kiểm tra chất lượng sản phẩm:} Trong các dây chuyền sản xuất công nghiệp, hệ thống phân loại ảnh tự động phát hiện sản phẩm lỗi, kiểm tra khuyết tật bề mặt với tốc độ và độ chính xác vượt trội so với kiểm tra thủ công.

\item \textbf{An ninh và giám sát:} Nhận dạng khuôn mặt, phân loại hành vi bất thường trong video giám sát, hỗ trợ lực lượng an ninh trong công tác phòng chống tội phạm.

\item \textbf{Viễn thám và môi trường:} Phân loại ảnh vệ tinh để giám sát sử dụng đất, theo dõi biến đổi rừng, dự báo thiên tai và đánh giá tác động môi trường.
\end{itemize}

\subsubsection{3.1.5. Các phương pháp tiếp cận bài toán phân loại ảnh}

\begin{quote}
Lịch sử phát triển của bài toán phân loại ảnh có thể chia thành hai giai đoạn chính:

\textbf{Giai đoạn 1 --- Phương pháp truyền thống (trước 2012):} Các phương pháp sử dụng kỹ thuật trích xuất đặc trưng thủ công (hand-crafted features) như SIFT (Scale-Invariant Feature Transform), HOG (Histogram of Oriented Gradients), hay LBP (Local Binary Pattern), kết hợp với các bộ phân lớp cổ điển như SVM (Support Vector Machine) hoặc KNN (K-Nearest Neighbors). Hạn chế chính của nhóm phương pháp này là đặc trưng được thiết kế bởi con người, phụ thuộc nhiều vào chuyên gia và không có khả năng tự thích ứng với dữ liệu mới.

\textbf{Giai đoạn 2 --- Phương pháp học sâu (từ 2012 đến nay):} Bước ngoặt lớn đến từ năm 2012 khi mô hình AlexNet của Krizhevsky và cộng sự giành chiến thắng cuộc thi ImageNet Large Scale Visual Recognition Challenge (ILSVRC) với khoảng cách vượt trội so với các phương pháp truyền thống. Kể từ đó, mạng nơ-ron tích chập (CNN) trở thành kiến trúc chủ đạo cho bài toán phân loại ảnh. Các kiến trúc CNN đặc biệt thành công bao gồm VGGNet (2014), GoogLeNet/Inception (2014), ResNet (2015), DenseNet (2017), và EfficientNet (2019). Ưu điểm cốt lõi của phương pháp học sâu là khả năng \textbf{tự động học các đặc trưng phân cấp} (hierarchical feature learning) trực tiếp từ dữ liệu thô, mà không cần sự can thiệp thủ công trong quá trình thiết kế đặc trưng.

Tuy nhiên, các mô hình CNN hiệu năng cao thường có kích thước rất lớn (hàng chục đến hàng trăm triệu tham số), đòi hỏi tài nguyên tính toán mạnh mẽ, không phù hợp với các thiết bị có tài nguyên hạn chế như điện thoại di động, thiết bị IoT hay thiết bị bay không người lái. Chính nhu cầu \textbf{cân bằng giữa hiệu năng phân loại và chi phí tính toán} đã thúc đẩy sự phát triển của các kiến trúc mạng CNN siêu nhẹ (light-weight CNNs) như MobileNets, ShuffleNets và \textbf{TickNets} --- là đối tượng nghiên cứu chính của tiểu luận này.
\end{quote}

\subsubsection{3.1.6. Vai trò của cơ chế chú ý trong bài toán phân loại ảnh}

\begin{quote}
Bên cạnh việc tối ưu kiến trúc mạng, một hướng nghiên cứu quan trọng khác nhằm nâng cao hiệu năng phân loại ảnh là tích hợp \textbf{cơ chế chú ý} (Attention Mechanism). Cơ chế chú ý mô phỏng khả năng tập trung chọn lọc của thị giác con người: thay vì xử lý đồng đều toàn bộ hình ảnh, mô hình học cách \textbf{nhấn mạnh các vùng và đặc trưng quan trọng} trong khi giảm thiểu ảnh hưởng của nhiễu nền và thông tin không liên quan.

Đặc biệt trong các bài toán phân loại ảnh thực tế --- nơi đối tượng cần nhận dạng thường bị lẫn trong nền phức tạp (ví dụ: vùng bệnh nhỏ trên lá cây giữa nền lá xanh, hay vùng viêm mờ trên ảnh X-quang phổi) --- cơ chế chú ý đóng vai trò quyết định trong việc giúp mô hình \textbf{định vị chính xác vùng đặc trưng mang tính phân biệt cao}, từ đó cải thiện đáng kể độ chính xác phân loại.

Chính từ bối cảnh trên, tiểu luận này đặt ra câu hỏi nghiên cứu: \textit{Làm thế nào để tích hợp hiệu quả cơ chế chú ý CBAM vào kiến trúc mạng siêu nhẹ TickNets nhằm nâng cao hiệu năng phân loại ảnh mà không làm bùng nổ chi phí tính toán?} Câu trả lời cho câu hỏi này sẽ được trình bày chi tiết trong phần đề xuất mô hình cải tiến phân cấp ở mục 3.2 tiếp theo.
\end{quote}

\subsection{3.2. ĐỀ XUẤT MÔ HÌNH TICKNETS KẾT HỢP CƠ CHẾ CHÚ Ý PHÂN CẤP (HIERARCHICAL ATTENTION):}\label{3-2-proposed-model}

Từ bối cảnh phân tích bài toán phân loại ảnh ở mục 3.1 và nền tảng lý thuyết về kiến trúc TickNets cùng cơ chế chú ý CBAM đã trình bày tại Chương 2, mục này đề xuất một kiến trúc cải tiến mang tên \textbf{TickNet-small CBAM-Hook}. Ý tưởng cốt lõi của đề xuất là thiết lập một \textbf{hệ thống chú ý phân cấp hai tầng} (hierarchical dual-level attention): giữ nguyên cơ chế chú ý kênh SE ở cấp độ cục bộ bên trong mỗi khối FR-PDP, đồng thời chèn mô-đun chú ý hỗn hợp CBAM ở cấp độ toàn cục tại các \textit{điểm nối chiến lược} (hooking points) giữa hai nhánh xương sống và trước lớp Global Average Pooling cuối mạng.

\subsubsection{3.2.1. Thiết kế kiến trúc đề xuất TickNet-small CBAM-Hook}\label{3-2-1-arch-design}

\paragraph{Động lực thiết kế.}

Như đã phân tích tại mục 1.2 và 1.4, kiến trúc TickNet-small nguyên bản sử dụng cấu trúc xương sống co giãn kênh hình dấu tích (Tick-shape Backbone) với các khối FR-PDP, trong đó mỗi khối tích hợp sẵn cơ chế chú ý kênh Squeeze-and-Excitation (SE) để tái hiệu chỉnh trọng số các kênh đặc trưng ở mức cục bộ. Tuy nhiên, mô hình gốc \textbf{hoàn toàn thiếu vắng cơ chế lọc đặc trưng theo không gian} --- tức khả năng xác định \textit{vùng nào} trên bản đồ đặc trưng chứa thông tin quan trọng nhất cho việc phân loại. Đây là một hạn chế đáng kể, đặc biệt khi xử lý các bộ dữ liệu thực tế có nền phức tạp như ảnh X-quang phổi (Chest-Xray-Pneumonia) hay ảnh bệnh lá cây (PlantVillage), nơi vùng bệnh lý thường chỉ chiếm một phần nhỏ của tổng diện tích ảnh và dễ bị lẫn trong nền.

Câu hỏi đặt ra là: \textit{Nên chèn cơ chế chú ý không gian CBAM vào vị trí nào trong kiến trúc TickNets để đạt hiệu quả tối ưu?} Nghiên cứu này xem xét và so sánh hai chiến lược chèn CBAM khác nhau:

\begin{itemize}
\item \textbf{Chiến lược chèn cục bộ (CBAM-Local):} Thay thế trực tiếp mô-đun SE bên trong mỗi khối FR-PDP bằng mô-đun CBAM. Đây là cách tiếp cận trực quan nhất, tương tự như cách CBAM được áp dụng trong các mạng CNN truyền thống (ví dụ: chèn sau mỗi convolutional block trong ResNet). Tuy nhiên, chiến lược này tiềm ẩn \textbf{xung đột kiến trúc} nghiêm trọng khi áp dụng vào các mạng siêu nhẹ sử dụng phép tích chập tách biệt chiều sâu, như sẽ được phân tích chi tiết tại mục 3.2.3.

\item \textbf{Chiến lược chèn phân cấp tại điểm nối (CBAM-Hook --- đề xuất cải tiến):} Giữ nguyên mô-đun SE cục bộ bên trong từng khối FR-PDP, đồng thời chèn CBAM ở \textit{bên ngoài} các khối tại những vị trí chiến lược dọc theo kiến trúc mạng. Chiến lược này tạo ra một hệ thống chú ý phân cấp hai tầng, trong đó mỗi tầng đảm nhận một vai trò khác biệt và bổ sung cho nhau.
\end{itemize}

\paragraph{Nguyên lý thiết kế kiến trúc CBAM-Hook.}

Kiến trúc TickNet-small CBAM-Hook được xây dựng dựa trên ba nguyên lý thiết kế cốt lõi:

\begin{enumerate}
\item \textbf{Nguyên lý phân tách chức năng (Separation of Concerns):} Cơ chế chú ý kênh SE ở mức cục bộ (nội khối) chịu trách nhiệm \textit{lọc kênh nào quan trọng} trong phạm vi một khối FR-PDP đơn lẻ. Cơ chế chú ý hỗn hợp CBAM ở mức toàn cục (ngoại khối) chịu trách nhiệm \textit{lọc vùng không gian nào quan trọng} trên toàn bộ bản đồ đặc trưng đã được tích lũy qua nhiều khối xử lý.

\item \textbf{Nguyên lý tránh chồng chéo đặc trưng (Feature Redundancy Avoidance):} Mô-đun SAM của CBAM sử dụng phép gộp toàn cục theo kênh (channel-wise pooling) để tạo bản đồ chú ý không gian 2D. Khi đặt SAM cùng mức cục bộ với phép tích chập chiều sâu (Depthwise Convolution) bên trong khối FR-PDP, hai phép toán này cùng hoạt động trên \textbf{cùng miền không gian cục bộ} với vùng tiếp nhận (receptive field) tương đồng, dẫn đến hiện tượng \textit{dư thừa thông tin không gian} (spatial information redundancy). Chi tiết phân tích toán học được trình bày tại mục 3.2.3.

\item \textbf{Nguyên lý chèn tại điểm nối chuyển pha (Phase Transition Hooking):} Kiến trúc Tick-shape Backbone của TickNets bao gồm hai pha rõ rệt: pha mở rộng kênh (expanding phase) và pha thu hẹp kênh (contracting phase). Điểm chuyển tiếp giữa hai pha chính là nơi đặc trưng đạt mức trừu tượng cao nhất với số kênh lớn nhất. Đây là vị trí lý tưởng để CBAM thực hiện việc \textit{tinh lọc toàn cục} --- chọn lọc cả kênh lẫn vùng không gian quan trọng nhất trước khi đặc trưng bước vào pha nén kênh.
\end{enumerate}

\paragraph{Kiến trúc chi tiết TickNet-small CBAM-Hook.}

Dựa trên ba nguyên lý trên, kiến trúc TickNet-small CBAM-Hook được thiết kế với ba điểm chèn CBAM chiến lược (hooking points), minh họa trong Hình \ref{fig:ticknet-cbam-hook-arch}:

\begin{figure}[htbp]
\centering
\includegraphics[width=0.85\textwidth]{./media/ticknet_cbam_hook_architecture.png}
\caption{Kiến trúc tổng thể mô hình đề xuất TickNet-small CBAM-Hook: các mô-đun CBAM được chèn tại ba điểm nối chiến lược (hooking points) dọc theo kiến trúc xương sống hình dấu tích}
\label{fig:ticknet-cbam-hook-arch}
\end{figure}

Cụ thể, ba vị trí chèn CBAM được xác định như sau:

\begin{itemize}
\item \textbf{Hooking Point 1 --- Điểm nối chuyển pha (giữa Expanding và Contracting Backbone):} Đây là vị trí mà số kênh đặc trưng đạt giá trị cực đại trong kiến trúc Tick-shape. Tại đây, đặc trưng đã đạt mức trừu tượng cao nhất sau khi đi qua toàn bộ pha mở rộng. CBAM tại vị trí này thực hiện hai chức năng quan trọng: (i) \textit{tái hiệu chỉnh kênh} thông qua Channel Attention Module (CAM) để chọn lọc những kênh đặc trưng giàu thông tin nhất trước khi bước vào pha nén, và (ii) \textit{lọc không gian} thông qua Spatial Attention Module (SAM) để loại bỏ nhiễu nền và tập trung vào vùng đối tượng chính.

\item \textbf{Hooking Point 2 --- Điểm nối cuối pha thu hẹp:} Sau khi đặc trưng đi qua toàn bộ pha thu hẹp kênh, số kênh giảm về mức nhỏ nhất. CBAM tại vị trí này đóng vai trò ``cổng kiểm soát cuối cùng'' (final gate), tinh chỉnh lại cả kênh lẫn không gian của bản đồ đặc trưng đã được nén trước khi truyền tới lớp phân loại.

\item \textbf{Hooking Point 3 --- Trước Global Average Pooling (GAP):} Đây là vị trí cuối cùng trước khi bản đồ đặc trưng 3D ($H \times W \times C$) bị nén thành vector 1D ($1 \times 1 \times C$) bởi phép GAP. CBAM tại đây thực hiện lần lọc không gian cuối cùng, đảm bảo rằng phép GAP chỉ tính trung bình trên các vùng không gian thực sự chứa đối tượng, không bị pha loãng bởi các vùng nền không liên quan.
\end{itemize}

Bảng \ref{tab:architecture-comparison} tóm tắt sự khác biệt kiến trúc giữa ba mô hình thực nghiệm.

\begin{table}[htbp]
\centering
\caption{So sánh kiến trúc ba mô hình thực nghiệm}
\label{tab:architecture-comparison}
\begin{tabular}{|l|c|c|c|}
\hline
\textbf{Thành phần kiến trúc} & \textbf{SE Baseline} & \textbf{CBAM-Local} & \textbf{CBAM-Hook} \\
\hline
Attention nội khối FR-PDP & SE & CBAM & SE (giữ nguyên) \\
\hline
Attention tại Hooking Point 1 & Không & Không & CBAM \\
\hline
Attention tại Hooking Point 2 & Không & Không & CBAM \\
\hline
Attention trước GAP & Không & Không & CBAM \\
\hline
Lọc kênh cục bộ & Có (SE) & Có (CAM) & Có (SE) \\
\hline
Lọc không gian cục bộ & Không & Có (SAM) & Không \\
\hline
Lọc không gian toàn cục & Không & Không & Có (SAM) \\
\hline
\end{tabular}
\end{table}

%%% ===== 3.2.2 =====
\subsubsection{3.2.2. Cơ chế chú ý phân cấp (Giữ SE nội khối FR-PDP, chèn CBAM ngoại khối tại các điểm nối Hooking points và cuối mạng)}\label{3-2-2-hierarchical-attention}

\paragraph{Khái niệm chú ý phân cấp.}

Thuật ngữ ``chú ý phân cấp'' (hierarchical attention) trong nghiên cứu này đề cập đến việc tổ chức các cơ chế chú ý theo \textbf{nhiều cấp độ trừu tượng} dọc theo chiều sâu của mạng, trong đó mỗi cấp độ đảm nhận một chức năng lọc đặc trưng khác biệt. Cụ thể, hệ thống phân cấp gồm hai tầng:

\begin{itemize}
\item \textbf{Tầng 1 --- Chú ý cục bộ nội khối (Intra-block Local Attention):} Cơ chế SE hoạt động \textit{bên trong} mỗi khối FR-PDP, thực hiện việc tái hiệu chỉnh trọng số kênh đặc trưng ở phạm vi cục bộ. SE nhận bản đồ đặc trưng đầu ra của phép tích chập trong khối, nén thông tin không gian qua Global Average Pooling thành vector $1 \times 1 \times C$, sau đó học hàm kích thích (excitation function) qua hai lớp fully connected để sinh ra vector trọng số kênh. Trọng số này được nhân ngược lại với bản đồ đặc trưng ban đầu, giúp khuếch đại các kênh chứa thông tin quan trọng và giảm thiểu các kênh nhiễu.

\item \textbf{Tầng 2 --- Chú ý toàn cục ngoại khối (Inter-block Global Attention):} Cơ chế CBAM hoạt động \textit{bên ngoài} các khối FR-PDP, tại các điểm nối chiến lược (hooking points). Khác với SE chỉ lọc theo kênh, CBAM thực hiện lọc tuần tự theo cả \textit{kênh} (qua CAM) và \textit{không gian} (qua SAM). Tại mỗi hooking point, CBAM nhận bản đồ đặc trưng đã được tích lũy và tinh lọc qua nhiều khối FR-PDP liên tiếp, từ đó thực hiện việc lọc ở \textbf{mức trừu tượng cao hơn} so với SE cục bộ.
\end{itemize}

\paragraph{Tại sao giữ nguyên SE nội khối thay vì thay thế bằng CBAM?}

Quyết định giữ nguyên SE bên trong khối FR-PDP thay vì thay thế bằng CBAM dựa trên ba luận cứ khoa học sau:

\textbf{Luận cứ 1 --- Tương thích kiến trúc với phép tích chập tách biệt:} Khối FR-PDP sử dụng chuỗi phép tính Pointwise Conv $\rightarrow$ Depthwise Conv $\rightarrow$ Pointwise Conv. Phép Depthwise Convolution (DWConv) xử lý mỗi kênh đặc trưng một cách \textit{độc lập} trên miền không gian cục bộ với kernel $k \times k$. Cơ chế SE chỉ hoạt động trên chiều kênh (channel-only) nên \textbf{hoàn toàn bổ sung} cho DWConv mà không có sự chồng chéo chức năng. Ngược lại, nếu thay SE bằng CBAM, thành phần SAM của CBAM sẽ cũng xử lý trên miền không gian cục bộ, tạo ra sự \textit{dư thừa chức năng} (functional redundancy) với DWConv (phân tích chi tiết tại mục 3.2.3).

\textbf{Luận cứ 2 --- Hiệu quả tham số:} Mô-đun SE cực kỳ nhẹ, chỉ bổ sung một lượng rất nhỏ tham số (hai lớp FC với tỉ lệ nén $r$). Trong khi đó, CBAM bổ sung thêm cả lớp convolution $7 \times 7$ cho SAM. Trong một kiến trúc siêu nhẹ như TickNets --- nơi tổng số tham số là yếu tố thiết kế quan trọng hàng đầu --- việc giữ SE trong mỗi khối giúp duy trì tính tối giản của mạng.

\textbf{Luận cứ 3 --- Phân chia tối ưu vai trò lọc đặc trưng:} Trong hệ thống phân cấp đề xuất, SE ở tầng cục bộ đóng vai trò ``bộ lọc kênh nhanh'' (fast channel filter) giúp mỗi khối FR-PDP tập trung vào các kênh đặc trưng quan trọng ngay trong quá trình xử lý nội bộ. CBAM ở tầng toàn cục đóng vai trò ``bộ lọc tổng hợp'' (aggregate filter) --- vừa tinh chỉnh kênh, vừa lọc không gian --- trên bản đồ đặc trưng đã trải qua nhiều bước xử lý, do đó mang tính tổng quát và có ý nghĩa ngữ nghĩa cao hơn.

\paragraph{Luồng xử lý đặc trưng qua kiến trúc phân cấp.}

Hình \ref{fig:three-model-comparison} minh họa sự khác biệt giữa ba chiến lược chèn cơ chế chú ý. Luồng xử lý đặc trưng trong kiến trúc CBAM-Hook đề xuất diễn ra như sau:

\begin{figure}[htbp]
\centering
\includegraphics[width=0.90\textwidth]{./media/hierarchical_attention_comparison.png}
\caption{So sánh ba chiến lược tích hợp cơ chế chú ý: (a) SE Baseline --- chỉ có SE nội khối, (b) CBAM-Local --- thay SE bằng CBAM trong khối, gây xung đột đặc trưng, (c) CBAM-Hook --- giữ SE nội khối, chèn CBAM phân cấp tại các điểm nối}
\label{fig:three-model-comparison}
\end{figure}

\textbf{Bước 1.} Ảnh đầu vào $\mathbf{X}_0 \in \mathbb{R}^{H \times W \times 3}$ đi qua lớp tích chập khởi tạo (stem convolution) để tạo bản đồ đặc trưng ban đầu.

\textbf{Bước 2.} Đặc trưng lần lượt đi qua các khối FR-PDP trong pha mở rộng kênh (Expanding Backbone). Trong mỗi khối, phép tính tuần tự PW $\rightarrow$ DW $\rightarrow$ PW được thực hiện kèm cơ chế SE nội khối và kết nối tắt toàn phần (full-residual skip connection). Gọi đầu ra sau khối thứ $i$ là $\mathbf{F}_i$:
\begin{equation}
\mathbf{F}_i = \text{FR-PDP}_i(\mathbf{F}_{i-1}) = \mathbf{F}_{i-1} + \text{SE}\big(\text{PW}_2(\text{DW}(\text{PW}_1(\mathbf{F}_{i-1})))\big)
\end{equation}

\textbf{Bước 3.} Tại Hooking Point 1 (sau khối cuối cùng của pha mở rộng), CBAM toàn cục được áp dụng:
\begin{equation}
\mathbf{F}'_{\text{hook1}} = \text{CBAM}(\mathbf{F}_N) = M_s\big(M_c(\mathbf{F}_N) \otimes \mathbf{F}_N\big) \otimes \big(M_c(\mathbf{F}_N) \otimes \mathbf{F}_N\big)
\end{equation}
trong đó $N$ là số khối của pha mở rộng, $M_c$ là bản đồ chú ý kênh, $M_s$ là bản đồ chú ý không gian, $\otimes$ là phép nhân theo phần tử.

\textbf{Bước 4.} Đặc trưng $\mathbf{F}'_{\text{hook1}}$ tiếp tục đi qua các khối FR-PDP trong pha thu hẹp kênh (Contracting Backbone), trải qua lọc SE cục bộ tương tự.

\textbf{Bước 5.} Tại Hooking Point 2 (sau khối cuối cùng của pha thu hẹp), CBAM toàn cục thứ hai được áp dụng tương tự.

\textbf{Bước 6.} Tại Hooking Point 3 (trước GAP), CBAM cuối cùng thực hiện lần lọc không gian quyết định, đảm bảo phép Global Average Pooling chỉ tổng hợp thông tin từ các vùng không gian mang tính phân biệt cao.

\textbf{Bước 7.} Lớp Global Average Pooling nén bản đồ đặc trưng 3D thành vector 1D, sau đó lớp Dense kết hợp hàm kích hoạt Softmax đưa ra phân phối xác suất trên các lớp đầu ra.

\paragraph{Ưu thế kỳ vọng của kiến trúc phân cấp.}

So với chiến lược SE Baseline (thiếu lọc không gian) và CBAM-Local (chồng chéo đặc trưng), kiến trúc CBAM-Hook đề xuất kỳ vọng mang lại các ưu thế sau:

\begin{itemize}
\item \textbf{Khả năng loại bỏ nhiễu nền vượt trội:} Nhờ ba lớp lọc không gian CBAM ở các điểm nối, mô hình có thể dần dần loại bỏ vùng nền không liên quan qua từng pha xử lý, đặc biệt hiệu quả trên các bộ dữ liệu thực tế có nền phức tạp.

\item \textbf{Khắc phục suy giảm hiệu năng trên ảnh nhỏ:} Trên các ảnh kích thước nhỏ (ví dụ: $28 \times 28$ của Fashion-MNIST), chiến lược CBAM-Local làm bản đồ đặc trưng nhanh chóng mất thông tin không gian do sự chồng chéo giữa DWConv và SAM. Chiến lược CBAM-Hook tránh được vấn đề này bằng cách chỉ áp dụng lọc không gian ở các điểm nối, nơi bản đồ đặc trưng đã đủ phong phú về ngữ nghĩa.

\item \textbf{Cân bằng giữa hiệu năng và chi phí tính toán:} Thay vì chèn CBAM vào mỗi khối (gây tăng tham số tuyến tính theo số khối), chỉ chèn tại ba điểm nối chiến lược giúp kiểm soát tổng lượng tham số bổ sung ở mức tối thiểu.
\end{itemize}

Các giả thuyết trên sẽ được kiểm chứng thông qua thực nghiệm đối chứng tại mục 3.3.

%%% ===== 3.2.3 =====
\subsubsection{3.2.3. Phân tích toán học giải quyết sự chồng chéo đặc trưng không gian giữa Depthwise Convolution và Spatial Attention Module (SAM)}\label{3-2-3-math-analysis}

Mục này trình bày phân tích toán học chặt chẽ để giải thích \textit{tại sao} chiến lược chèn CBAM cục bộ bên trong khối FR-PDP (CBAM-Local) có thể dẫn đến suy giảm hiệu năng, từ đó củng cố cơ sở khoa học cho chiến lược chèn phân cấp CBAM-Hook.

\paragraph{Phân tích vùng tiếp nhận (Receptive Field Analysis).}

Xét bản đồ đặc trưng trung gian $\mathbf{F} \in \mathbb{R}^{H \times W \times C}$ tại vị trí ngay sau phép Depthwise Convolution bên trong một khối FR-PDP. Tại mỗi vị trí không gian $(i, j)$, giá trị đặc trưng trên kênh thứ $c$ được tính bởi:

\begin{equation}
\mathbf{F}(i, j, c) = \sum_{m=-\lfloor k/2 \rfloor}^{\lfloor k/2 \rfloor} \sum_{n=-\lfloor k/2 \rfloor}^{\lfloor k/2 \rfloor} \mathbf{W}_c(m, n) \cdot \mathbf{X}(i+m, j+n, c) + b_c
\label{eq:dwconv}
\end{equation}

trong đó $\mathbf{W}_c \in \mathbb{R}^{k \times k}$ là kernel riêng cho kênh $c$, $k$ là kích thước kernel (thường $k = 3$ trong TickNets), và $\mathbf{X}$ là đầu vào. Điểm mấu chốt là: phép DWConv đã mã hóa thông tin không gian cục bộ trong vùng tiếp nhận $k \times k$ xung quanh mỗi vị trí $(i, j)$ vào giá trị $\mathbf{F}(i, j, c)$.

Bây giờ, xét hoạt động của SAM khi được đặt ngay sau DWConv trong khối FR-PDP. SAM tạo bản đồ chú ý không gian qua ba bước:

\textbf{Bước 1 --- Gộp kênh (Channel-wise pooling):} SAM áp dụng đồng thời Average Pooling và Max Pooling \textit{dọc theo chiều kênh} tại mỗi vị trí không gian:
\begin{equation}
\mathbf{F}_{\text{avg}}^s(i,j) = \frac{1}{C}\sum_{c=1}^{C} \mathbf{F}(i,j,c), \quad
\mathbf{F}_{\text{max}}^s(i,j) = \max_{c \in \{1,\ldots,C\}} \mathbf{F}(i,j,c)
\label{eq:sam-pool}
\end{equation}

\textbf{Bước 2 --- Tích chập $7 \times 7$:} Hai bản đồ $\mathbf{F}_{\text{avg}}^s$ và $\mathbf{F}_{\text{max}}^s$ được nối (concatenate) và đưa qua lớp tích chập với kernel $7 \times 7$:
\begin{equation}
M_s(i,j) = \sigma\Bigg(\sum_{m=-3}^{3}\sum_{n=-3}^{3} \mathbf{W}^s(m,n) \cdot [\mathbf{F}_{\text{avg}}^s(i+m,j+n);\; \mathbf{F}_{\text{max}}^s(i+m,j+n)] + b^s\Bigg)
\label{eq:sam-conv}
\end{equation}
trong đó $\sigma$ là hàm sigmoid, $\mathbf{W}^s \in \mathbb{R}^{7 \times 7 \times 2}$ là kernel của SAM.

\paragraph{Hiện tượng chồng chéo vùng tiếp nhận.}

Thay phương trình (\ref{eq:dwconv}) vào (\ref{eq:sam-pool}), ta thấy rằng giá trị $\mathbf{F}_{\text{avg}}^s(i,j)$ thực chất là trung bình có trọng số của các giá trị pixel trong vùng $k \times k$ quanh $(i,j)$ trên \textit{tất cả} $C$ kênh:

\begin{equation}
\mathbf{F}_{\text{avg}}^s(i,j) = \frac{1}{C}\sum_{c=1}^{C}\Bigg[\sum_{m,n} \mathbf{W}_c(m,n) \cdot \mathbf{X}(i+m,j+n,c) + b_c\Bigg]
\label{eq:overlap}
\end{equation}

Sau đó, phép tích chập $7 \times 7$ của SAM (phương trình \ref{eq:sam-conv}) lại tổng hợp thông tin từ vùng $7 \times 7$ xung quanh $(i,j)$ trên bản đồ $\mathbf{F}_{\text{avg}}^s$. Do đó, \textbf{vùng tiếp nhận hiệu dụng} (effective receptive field) của SAM khi đặt sau DWConv là:

\begin{equation}
\text{ERF}_{\text{SAM}} = (k - 1) + (7 - 1) + 1 = k + 5
\label{eq:erf}
\end{equation}

Với $k = 3$ (kernel DWConv điển hình trong TickNets), $\text{ERF}_{\text{SAM}} = 8$ pixel. Điều này có nghĩa là SAM cục bộ đang cố gắng học bản đồ chú ý không gian dựa trên thông tin trong vùng tiếp nhận chỉ $8 \times 8$ pixel --- thực chất đây vẫn là thông tin \textbf{cục bộ}, không mang tính toàn cục.

\paragraph{Phân tích dư thừa thông tin trên ảnh kích thước nhỏ.}

Vấn đề trở nên nghiêm trọng hơn trên các ảnh có kích thước nhỏ. Xét tập dữ liệu Fashion-MNIST với ảnh gốc $28 \times 28$ pixel. Sau một vài khối FR-PDP với stride $> 1$ hoặc pooling, kích thước bản đồ đặc trưng giảm nhanh chóng. Giả sử tại khối FR-PDP thứ $l$, bản đồ đặc trưng có kích thước $H_l \times W_l$. Khi $H_l$ hoặc $W_l$ giảm xuống gần bằng hoặc nhỏ hơn $\text{ERF}_{\text{SAM}}$, ta có:

\begin{equation}
\text{Tỉ lệ phủ} = \frac{\text{ERF}_{\text{SAM}}^2}{H_l \times W_l} \rightarrow 1 \quad \text{khi } H_l, W_l \rightarrow k + 5
\label{eq:coverage}
\end{equation}

Khi tỉ lệ phủ tiến gần 1, bản đồ chú ý không gian $M_s$ trở nên \textbf{gần như đồng nhất} trên toàn bộ bản đồ đặc trưng --- tức SAM mất khả năng phân biệt vùng không gian nào quan trọng hơn vùng nào. Trong trường hợp này, SAM không những không cung cấp thông tin bổ sung mà còn đưa thêm \textit{nhiễu gradient} vào quá trình huấn luyện do phải tối ưu các tham số $\mathbf{W}^s$ trên tín hiệu gần đồng nhất.

\paragraph{Tại sao CBAM-Hook giải quyết được vấn đề?}

Trong kiến trúc CBAM-Hook, mô-đun SAM không được đặt bên trong từng khối FR-PDP mà được đặt tại các hooking points --- nơi bản đồ đặc trưng đã trải qua nhiều khối xử lý liên tiếp. Tại các vị trí này, mỗi pixel trên bản đồ đặc trưng đã tích lũy thông tin từ một vùng tiếp nhận \textit{lớn hơn nhiều} trên ảnh gốc nhờ hiệu ứng chồng tầng (cascading effect) của nhiều lớp tích chập:

\begin{equation}
\text{ERF}_{\text{tích lũy}}^{(L)} = 1 + \sum_{l=1}^{L} (k_l - 1) \cdot \prod_{j=l+1}^{L} s_j
\label{eq:cumulative-erf}
\end{equation}

trong đó $L$ là số lớp tích chập đã đi qua, $k_l$ là kích thước kernel tại lớp $l$, $s_j$ là stride tại lớp $j$. Với $L$ đủ lớn (tại hooking points), $\text{ERF}_{\text{tích lũy}}$ có thể bao phủ phần lớn hoặc toàn bộ ảnh đầu vào. Khi đó, phép gộp kênh của SAM (phương trình \ref{eq:sam-pool}) thực hiện trên bản đồ đặc trưng mà mỗi vị trí đã chứa thông tin \textbf{ngữ nghĩa toàn cục}, giúp SAM thực sự có khả năng phân biệt ``đâu là vùng đối tượng'' và ``đâu là vùng nền'' một cách có ý nghĩa.

\paragraph{Tóm tắt luận điểm toán học.}

Phân tích trên cho thấy hai kết luận quan trọng:

\begin{enumerate}
\item \textbf{CBAM-Local gây dư thừa:} Khi SAM được đặt ngay sau DWConv bên trong khối FR-PDP, vùng tiếp nhận hiệu dụng của SAM bị giới hạn ở mức cục bộ ($k + 5$ pixel), chồng chéo chức năng với DWConv. Trên ảnh nhỏ, bản đồ chú ý không gian trở nên đồng nhất và vô nghĩa, gây suy giảm hiệu năng.

\item \textbf{CBAM-Hook tối ưu:} Khi SAM được đặt tại các hooking points sau nhiều khối FR-PDP, vùng tiếp nhận tích lũy đủ lớn để SAM hoạt động trên thông tin ngữ nghĩa toàn cục. Đồng thời, việc giữ nguyên SE nội khối đảm bảo lọc kênh cục bộ vẫn hoạt động hiệu quả mà không có xung đột chức năng.
\end{enumerate}

Các kết luận lý thuyết này sẽ được kiểm chứng bằng thực nghiệm đối chứng trên bốn bộ dữ liệu ở mục 3.3, đặc biệt qua phân tích hiện tượng suy giảm hiệu năng cục bộ (mục 3.3.2.4) và trực quan hóa Grad-CAM (mục 3.3.2.5).

\subsection{3.3. HIỆN THỰC GIẢI PHÁP VÀ THỰC NGHIỆM:}\label{hiux1ec7n-thux1ef1c-giux1ea3i-phuxe1p}

\subsubsection{3.3.1. Môi trường thực nghiệm và xây dựng mô hình:}\label{muxf4i-trux1b0ux1eddng-thux1ef1c-nghiux1ec7m-vuxe0-xuxe2y-dux1ef1ng-muxf4-huxecnh}

Quá trình xây dựng mô hình và huấn luyện bao gồm các giai đoạn:

\pandocbounded{\includegraphics[keepaspectratio]{./media/image34.png}}

\begin{quote}
\emph{Hình 3.3 Quá trình xây dựng và huấn luyện mô hình}
\end{quote}

\paragraph{3.3.1.1. Môi trường thực nghiệm (Nền tảng Kaggle, GPU NVIDIA T4/P100, TensorFlow và Keras)}\label{3-3-1-1-environment}
\begin{quote}
Ngôn ngữ lập trình để thực hiện quá trình thực nghiệm là Python với các
thư viện hỗ trợ như Keras, TensorFlow, Pytorch, Numpy,
Scikit-learn\ldots{}

Quá trình xây dựng và huấn luyện mô hình được thực hiện trên nền tảng điện
toán đám mây Kaggle Notebooks chạy hệ điều hành Linux, với dung lượng ổ
đĩa khả dụng khoảng 73.0 GB, 13GB--30GB RAM hệ thống và bộ tăng tốc card
đồ họa GPU NVIDIA T4 (15GB VRAM) hoặc GPU P100 (16GB VRAM).
\end{quote}

\paragraph{3.3.1.2. Mô tả 4 bộ dữ liệu thực nghiệm:}\label{3-3-1-2-datasets}

\subparagraph{3.3.1.2.1. Nhóm dữ liệu chuẩn học thuật: Fashion-MNIST và CIFAR-10}\label{3-3-1-2-1-standard}
\begin{quote}
Chuẩn bị tập dữ liệu cho thực nghiệm bao gồm tập huấn luyện và tập kiểm
tra. Các tập dữ liệu được chọn để thực hành nghiên cứu bao gồm
Fashion-MNIST và CIFAR-10.
\end{quote}

\begin{itemize}
\item
  \textbf{Fashion-MNIST:} là tập dữ liệu chứa các hình ảnh thuộc 10 lớp,
  mỗi lớp có 6,000 ảnh trong tập huấn luyện và 1,000 ảnh trong tập kiểm
  tra. Do đó, tập huấn luyện và tập kiểm tra sẽ chứa tổng cộng lần lượt
  60,000 và 10,000 ảnh.
\end{itemize}

\begin{quote}
Các ảnh trong Fashion-MNIST tương ứng với các lớp: Áo phông/T-shirt
(T-shirt/top), Quần (Trouser), Áo len (Pullover), Váy (Dress), Áo khoác
(Coat), Dép (Sandal), Áo sơ mi (Shirt), Giày thể thao (Sneaker), Túi
xách (Bag), Bốt/giày cổ cao (Ankle boot). Mỗi hình ảnh có kích thước
28x28 pixel và là ảnh xám (grayscale).
\end{quote}

\paragraph{Bảng 3.1 Bảng mô tả dữ liệu tập dữ liệu
MNIST-Fashion}\label{bux1ea3ng-3.1-bux1ea3ng-muxf4-tux1ea3-dux1eef-liux1ec7u-tux1eadp-dux1eef-liux1ec7u-mnist-fashion}

{\def\LTcaptype{none} % do not increment counter
\begin{longtable}[]{@{}
  >{\raggedright\arraybackslash}p{(\linewidth - 2\tabcolsep) * \real{0.2068}}
  >{\raggedright\arraybackslash}p{(\linewidth - 2\tabcolsep) * \real{0.6070}}@{}}
\toprule\noalign{}
\begin{minipage}[b]{\linewidth}\raggedright
\textbf{Thuộc tính}
\end{minipage} & \begin{minipage}[b]{\linewidth}\raggedright
\begin{quote}
\textbf{Mô tả}
\end{quote}
\end{minipage} \\
\midrule\noalign{}
\endhead
\bottomrule\noalign{}
\endlastfoot
\textbf{Số lượng ảnh} & \begin{minipage}[t]{\linewidth}\raggedright
\begin{quote}
70,000
\end{quote}
\end{minipage} \\
\textbf{Số lớp} & \begin{minipage}[t]{\linewidth}\raggedright
\begin{quote}
10
\end{quote}
\end{minipage} \\
\textbf{Kích thước ảnh} & \begin{minipage}[t]{\linewidth}\raggedright
\begin{quote}
28x28 pixels
\end{quote}
\end{minipage} \\
\textbf{Định dạng} & \begin{minipage}[t]{\linewidth}\raggedright
\begin{quote}
Grayscale (ảnh đen trắng)
\end{quote}
\end{minipage} \\
\textbf{Số lượng tập}

\textbf{huấn luyện} & \begin{minipage}[t]{\linewidth}\raggedright
\begin{quote}
60,000
\end{quote}
\end{minipage} \\
\textbf{Số lượng tập}

\textbf{kiểm tra} & \begin{minipage}[t]{\linewidth}\raggedright
\begin{quote}
10,000
\end{quote}
\end{minipage} \\
\textbf{Tên các lớp} & \begin{minipage}[t]{\linewidth}\raggedright
\begin{quote}
0: T-shirt/top, 1: Trouser, 2: Pullover, 3: Dress, 4:

Coat, 5: Sandal, 6: Shirt, 7: Sneaker, 8: Bag, 9: Ankle boot.
\end{quote}
\end{minipage} \\
\end{longtable}
}

\begin{itemize}
\item
  \textbf{CIFAR-10}: là tập dữ liệu bao gồm 60,000 hình ảnh có kích
  thước 32x32 pixel và là ảnh màu (RGB) được chia thành 10 lớp, với
  6,000 hình ảnh mỗi lớp. Có 50,000 hình ảnh trong tập huấn luyện và
  10,000 hình ảnh trong tập kiểm tra. Các ảnh trong CIFAR-10 tương ứng
  với các lớp: Máy bay (Airplane), Xe hơi (Automobile), Chim (Bird), Mèo
  (Cat), Hươu (Deer), Chó (Dog), Ếch (Frog), Ngựa (Horse), Tàu thủy
  (Ship), Xe tải (Truck).
\end{itemize}

\paragraph{Bảng 3.2 Bảng mô tả dữ liệu tập dữ liệu
Cifar-10}\label{bux1ea3ng-3.2-bux1ea3ng-muxf4-tux1ea3-dux1eef-liux1ec7u-tux1eadp-dux1eef-liux1ec7u-cifar-10}

{\def\LTcaptype{none} % do not increment counter
\begin{longtable}[]{@{}
  >{\raggedright\arraybackslash}p{(\linewidth - 2\tabcolsep) * \real{0.2069}}
  >{\raggedright\arraybackslash}p{(\linewidth - 2\tabcolsep) * \real{0.6071}}@{}}
\toprule\noalign{}
\begin{minipage}[b]{\linewidth}\raggedright
\textbf{Thuộc tính}
\end{minipage} & \begin{minipage}[b]{\linewidth}\raggedright
\begin{quote}
\textbf{Mô tả}
\end{quote}
\end{minipage} \\
\midrule\noalign{}
\endhead
\bottomrule\noalign{}
\endlastfoot
\textbf{Số lượng ảnh} & \begin{minipage}[t]{\linewidth}\raggedright
\begin{quote}
60,000
\end{quote}
\end{minipage} \\
\textbf{Số lớp} & \begin{minipage}[t]{\linewidth}\raggedright
\begin{quote}
10
\end{quote}
\end{minipage} \\
\textbf{Kích thước ảnh} & \begin{minipage}[t]{\linewidth}\raggedright
\begin{quote}
32x32 pixels
\end{quote}
\end{minipage} \\
\textbf{Định dạng} & \begin{minipage}[t]{\linewidth}\raggedright
\begin{quote}
RGB (ảnh màu).
\end{quote}
\end{minipage} \\
\textbf{Số lượng tập}

\textbf{huấn luyện} & \begin{minipage}[t]{\linewidth}\raggedright
\begin{quote}
50,000
\end{quote}
\end{minipage} \\
\textbf{Số lượng tập}

\textbf{kiểm tra} & \begin{minipage}[t]{\linewidth}\raggedright
\begin{quote}
10,000
\end{quote}
\end{minipage} \\
\textbf{Tên các lớp} & \begin{minipage}[t]{\linewidth}\raggedright
\begin{quote}
0: Airplane, 1: Automobile, 2: Bird, 3: Cat, 4: Deer,

5: Dog, 6: Frog, 7: Horse, 8: Ship, 9: Truck.
\end{quote}
\end{minipage} \\
\end{longtable}
}

\begin{quote}
Cả hai tập dữ liệu đều có thể tải trược tiếp từ trang web kaggle hoặc
thông qua thư viện Tensorflow.

\includegraphics[width=5.17775in,height=0.34375in]{./media/image38.jpeg}
\end{quote}

\includegraphics[width=5.07905in,height=0.375in]{./media/image39.jpeg}

\begin{quote}
\emph{Hình 3.4 Tải tập dữ liệu}
\end{quote}

\subparagraph{3.3.1.2.2. Nhóm dữ liệu ứng dụng thực tiễn: PlantVillage (bệnh lá cây) và Chest-Xray-Pneumonia (viêm phổi)}\label{3-3-1-2-2-practical}
\begin{quote}
\emph{[Nội dung phần này đang được cập nhật/viết tiếp...]}
\end{quote}

\paragraph{3.3.1.3. Các phương pháp Tiền xử lý dữ liệu và Tăng cường dữ liệu (Data Augmentation) cho từng nhóm ảnh}\label{3-3-1-3-preprocessing}

\begin{quote}
\emph{Chuẩn hóa dữ liệu (Normalization):} Dữ liệu hình ảnh thường được
lưu trữ dưới dạng số nguyên từ 0 đến 255, tương ứng với cường độ pixel.
Chuẩn hóa dữ liệu chuyển đổi các giá trị này về khoảng {[}0, 1{]} bằng
cách chia cho

255. Quá trình này giúp các thuật toán học máy hội tụ nhanh hơn và đạt
hiệu quả tốt hơn.

\emph{One hot coding:} Quá trình này chuyển đổi các nhãn lớp từ dạng số
nguyên sang dạng vector nhị phân. Mỗi nhãn số nguyên được chuyển thành
một vector nhị phân có độ dài bằng số lớp (ở đây là 10), với một phần tử
là 1 và các phần tử còn lại là 0. Ví dụ, nhãn 3 sẽ được chuyển thành
{[}0, 0, 0, 1, 0, 0, 0, 0, 0, 0{]}. Giúp các mô hình học sâu (như mạng
nơ-ron) có thể xử lý các nhãn này một cách hiệu quả hơn.
\end{quote}

\paragraph{3.3.1.4. Hiện thực hóa 3 mô hình đối chứng thực nghiệm:}\label{3-3-1-4-models}

\subparagraph{3.3.1.4.1. Mô hình 1 (SE Baseline): TickNet-small nguyên bản sử dụng chú ý SE trong khối}\label{3-3-1-4-1-baseline}
\begin{quote}
\emph{[Nội dung phần này đang được cập nhật/viết tiếp...]}
\end{quote}

\subparagraph{3.3.1.4.2. Mô hình 2 (CBAM-Local): Chèn trực tiếp CBAM thay SE cục bộ trong khối FR-PDP}\label{3-3-1-4-2-local}
\begin{quote}
\emph{[Nội dung phần này đang được cập nhật/viết tiếp...]}
\end{quote}

\subparagraph{3.3.1.4.3. Mô hình 3 (CBAM-Hook - Đề xuất cải tiến): Chèn CBAM phân cấp tại các điểm nối và cuối mạng}\label{3-3-1-4-3-hook}
\begin{quote}
\emph{[Nội dung phần này đang được cập nhật/viết tiếp...]}
\end{quote}

\textbf{Xây dựng và cấu trúc mạng CNN thực nghiệm:}

\begin{quote}
\emph{\textbf{Transfer learning:}} là một kỹ thuật trong học máy, trong
đó một mô hình được phát triển cho một nhiệm vụ cụ thể được tái sử dụng
làm điểm khởi đầu cho một mô hình trên một nhiệm vụ thứ hai. Thay vì bắt
đầu từ đầu (từ dữ liệu gốc và mô hình mới), học chuyển giao bắt đầu từ
các đặc trưng đã học từ một mô hình đã được huấn luyện trước đó.
{[}12{]}

\emph{\textbf{Fine tuning:}} là một phương pháp của transfer learning,
lấy 1 pre-trained model như VGG, Resnet \ldots{} tận dụng 1 phần hoặc
toàn bộ các lớp, thêm/sửa/xoá 1 vài lớp để tạo ra 1 model mới. Thường
các lớp đầu của model được freeze (đóng băng) lại - tức weight các layer
này sẽ không bị thay đổi giá trị trong quá trình train. Lý do bởi các
lớp này đã có khả năng trích xuất thông tin mức trừu tượng thấp, khả
năng này được học từ quá trình training trước đó. Ta freeze lại để tận
dụng được khả năng này và giúp việc train diễn ra nhanh hơn (model chỉ
phải cập nhật trọng số ở các lớp cao). {[}13{]}

Đối với tập dữ liệu MSNIST-Fashion sẽ được thay đổi kích thước ảnh thành
64x64x3 để phù hợp với việc áp dụng fine tuning vào mô hình. Tải
pre-trained model VGG16 từ thư viện Tensorflow và loại bỏ lớp FC. Thêm
vào các lớp phù hợp với tập dữ liệu. Đóng băng các lớp đầu của mô hình
VGG 16 và hủy đóng băng ở một vài lớp cuối cùng.

\includegraphics[width=4.84371in,height=7.75625in]{./media/image40.png}

\emph{Hình 3.5 Mô hình CNN cho tập dữ liệu MNIST-Fashion}

\includegraphics[width=4.64069in,height=8.10625in]{./media/image41.png}

\emph{Hình 3.6 Mô hình CNN cho tập dữ liệu Cifar-10}

\emph{Đối với tập dữ liệu MNIST-Fashion:} ngoài pre-trained model VGG16
sẽ thêm một lớp maxpool với cửa sổ pooling có kích thước là 2x2 và
stride là (2,2). Và thêm lớp Flattten để làm phẳng, tiếp theo là một lớp
Dropout với tỉ lệ là 0.5. Tục là các lớp fully connected với số nơ-ron
cho các lớp lần lượt là 1000,800,400,200,100 với hàm kích hoạt là ReLU.
Cuối cùng là lớp đầu ra với 10 nơ-ron tương ứng với số nhãn của tập dữ
liệu với một hàm kích hoạt softmax để phân bổ xác xuất.

\emph{Đối với tập dữ liệu Cifar-10:} ngoài pre-trained model VGG16 sẽ
thêm lớp Flattten để làm phẳng, tiếp theo là một lớp Dropout với tỉ lệ
là 0.5. Tục là các lớp fully connected với số nơ-ron cho các lớp lần
lượt là 516, 256, 128, 64 với hàm kích hoạt là ReLU và 2 lớp Drop out
với tỉ lệ 0.5. Cuối cùng là lớp đầu ra với 10 nơ-ron tương ứng với số
nhãn của tập dữ liệu với một hàm kích hoạt softmax để phân bổ xác xuất.
\end{quote}

\textbf{Áp dụng cơ chế chú ý (Attention Mechanism) vào mô hình CNN:}

\begin{quote}
\emph{Đối với tập dữ liệu MNIST-Fashion:} Một cơ chế CBAM sẽ được thêm
vào sau lớp Convolution cuối cùng (block5\_conv3) của pre-trained model
Với tỉ lệ giảm r = 12.

\emph{Đối với tập dữ liệu Cifar-10:} Cơ chế CBAM sẽ được thêm vào
Convolution cuối của block3, block4 và block5 (block3\_conv3,
block4\_conv3, block5\_conv3) của pre-trained model cùng với hàm kích
hoạt phi tuyến ReLU và các lớp BatchNormalization. Với tỉ lệ giảm r =2.
\end{quote}

\paragraph{3.3.1.5. Cấu hình quá trình huấn luyện và tối ưu hóa:}\label{3-3-1-5-training-config}

\subparagraph{3.3.1.5.1. Bộ tối ưu SGD và các hàm Loss tương ứng}\label{3-3-1-5-1-sgd-loss}

\begin{quote}
Cả 2 tập dữ liệu cho mô hình CNN đều được huấn luyện với hàm loss là
``categorical\_crossentropy'', bộ tối ưu hóa là SGD, metric là
``accuracy'', Validation\_plit là 0.1 với 10 epochs và bath\_size là 32.
\end{quote}

\subparagraph{3.3.1.5.2. Hàm gọi lại tự động giảm tốc độ học (ReduceLROnPlateau) và dừng sớm (EarlyStopping)}\label{3-3-1-5-2-callbacks}
\begin{quote}
\emph{[Nội dung phần này đang được cập nhật/viết tiếp...]}
\end{quote}

\subsubsection{3.3.2. Kết quả thực nghiệm, phân tích và đánh giá:}\label{kux1ebft-quux1ea3-thux1ef1c-nghiux1ec7m-vuxe0-ux111uxe1nh-giuxe1}

\paragraph{3.3.2.1. Các độ đo hiệu năng chuẩn mực (Accuracy, Precision, Recall, F1-score) và cấu trúc Confusion Matrix}\label{3-3-2-1-metrics}

\begin{quote}
Confusion matrix là một công cụ đánh giá hiệu suất của các mô hình phân
loại trong học máy. Nó cung cấp một cái nhìn chi tiết về cách mà mô hình
phân loại các

mẫu trong tập dữ liệu kiểm tra (hoặc tập dữ liệu kiểm chứng). Ma trận
này chứa các giá trị tổng hợp cho các dự đoán đúng và sai của mô hình,
được phân loại dựa trên các lớp thực tế và dự đoán.
\end{quote}

\includegraphics[width=4.6131in,height=2.55667in]{./media/image42.jpeg}

\begin{quote}
\emph{Hình 3.7 Cấu trúc Confusion matrix}

Trong đó:
\end{quote}

\begin{itemize}
\item
  True postitive (TP): các điểm dữ liệu dự đoán là 1 và có nhãn là 1.
\item
  True negative (TN): các điểm dữ liệu dự đoán là 0 và có nhãn là 0.
\item
  False positive (FP): các điểm dữ liệu dự đoán là 1 nhưng có nhãn là 0.
\item
  False negative (FN): các điểm dữ liệu dự đoán là 0 nhưng có nhãn là 1.
\end{itemize}

\begin{quote}
Confusion Matrix cung cấp thông tin quan trọng để tính toán các chỉ số
đánh giá hiệu suất của mô hình phân loại như:
\end{quote}

\begin{itemize}
\item
  Accuracy: Cách đơn giản nhất để đánh giá một mô hình phân lớp đó là sử
  dụng độ chính xác (Accuracy). Ý tưởng đơn giản là tỷ lệ giữa các mẫu
  dự đoán đúng trên tổng số mẫu của dữ liệu kiểm thử.
\end{itemize}

\begin{quote}
Accuracy = \textsuperscript{TP+TN}

TP+TN+FP+FN
\end{quote}

\begin{itemize}
\item
  Precision: Đo lường tỷ lệ các dự đoán tích cực thực sự trong số tất cả
  các dự đoán tích cực do mô hình đưa ra.
\end{itemize}

\begin{quote}
Precision = \textsuperscript{TP}

TP+F
\end{quote}

\begin{itemize}
\item
  Recall: Đo lường tỷ lệ dự đoán tích cực thực sự trong số tất cả các
  trường hợp tích cực thực tế.
\end{itemize}

\begin{quote}
Recall = \textsuperscript{TP}

TP+FN
\end{quote}

\begin{itemize}
\item
  F1 Score: F1 Score là giá trị trung bình của precision và recall.
  F1-score có giá trị nằm trong nửa khoảng (𝟎, 𝟏{]}. F1-score càng cao
  đồng nghĩa với bộ phân lớp càng tốt. Khi cả recall và precision đều
  bằng 1 thì F1-score = 1. Khi cả recall và precision đều thấp thì
  F1-score tiến về 0.
\end{itemize}

\begin{quote}
F1 Score = 2× precision×recall

precision+recal

Để đánh giá hiệu quả của các mô hình học viên sẽ dựa vào các chỉ số như
accuracy, precision, recall và F1 score trên cả hai tập dữ liệu.
\end{quote}

\paragraph{3.3.2.2. Kết quả thực nghiệm và biểu đồ Loss/Accuracy trên tập dữ liệu chuẩn (Fashion-MNIST, CIFAR-10)}\label{3-3-2-2-standard-results}

\begin{quote}
\includegraphics[width=4.94816in,height=4.00937in]{./media/image43.png}

\emph{Hình 3.8 Biểu đồ loss và accuracy của mô hình CNN có áp dụng
Attention và không áp dụng với tập dữ liệu MNIST-Fashion}

Nhìn chung cả hai mô hình đều có xu hướng giảm loss tương tự nhau.
Training loss và validation loss của cả hai mô hình đều giảm dần và giữ
xu hướng giảm gần như song song sau khoảng epoch 2. Mặc dù cả hai mô
hình có xu hướng rất giống nhau, nhưng mô hình với Attention có accuracy
ổn định và cao hơn một chút so với mô hình không có Attention.

\includegraphics[width=5.20692in,height=3.95in]{./media/image44.jpeg}

\emph{Hình 3.9 Biểu đồ loss và accuracy của mô hình CNN có áp dụng
Attention và không áp dụng với tập dữ liệu Cifar-10}

Cả hai mô hình đều cho thấy giảm training và validation loss qua các
epoch. Mô hình có attention dường như có đường cong validation loss mượt
hơn một chút, cho thấy quá trình học ổn định hơn. Về Accuracy thì mô
hình có attention có sự tăng trưởng training accuracy ổn định hơn.
Validation accuracy của mô hình có attention cũng ổn định hơn so với mô
hình không có attention, mô hình không có attention có những dao động rõ
rệt hơn. Qua đó có thể thấy việc thêm các cơ chế attention vào CNN dường
như giúp ổn định quá trình huấn luyện và có thể cải thiện khả năng của
mô hình trong việc tổng quát hóa với dữ liệu chưa thấy trước đó.
\end{quote}

\paragraph{Bảng 3.3 Bảng so sánh kết quả giữa các giải thuật trên tập dữ
liệu
MNIST-Fashion}\label{bux1ea3ng-3.3-bux1ea3ng-so-suxe1nh-kux1ebft-quux1ea3-giux1eefa-cuxe1c-giux1ea3i-thuux1eadt-truxean-tux1eadp-dux1eef-liux1ec7u-mnist-fashion}

{\def\LTcaptype{none} % do not increment counter
\begin{longtable}[]{@{}
  >{\raggedright\arraybackslash}p{(\linewidth - 8\tabcolsep) * \real{0.3863}}
  >{\raggedright\arraybackslash}p{(\linewidth - 8\tabcolsep) * \real{0.1339}}
  >{\raggedright\arraybackslash}p{(\linewidth - 8\tabcolsep) * \real{0.1337}}
  >{\raggedright\arraybackslash}p{(\linewidth - 8\tabcolsep) * \real{0.1489}}
  >{\raggedright\arraybackslash}p{(\linewidth - 8\tabcolsep) * \real{0.1181}}@{}}
\toprule\noalign{}
\begin{minipage}[b]{\linewidth}\centering
\textbf{Algorithm}
\end{minipage} & \begin{minipage}[b]{\linewidth}\raggedright
\begin{quote}
\textbf{Accuracy}
\end{quote}
\end{minipage} & \begin{minipage}[b]{\linewidth}\raggedright
\begin{quote}
\textbf{Precision}
\end{quote}
\end{minipage} & \begin{minipage}[b]{\linewidth}\raggedright
\begin{quote}
\textbf{Recall}
\end{quote}
\end{minipage} & \begin{minipage}[b]{\linewidth}\centering
\textbf{F1}
\end{minipage} \\
\midrule\noalign{}
\endhead
\bottomrule\noalign{}
\endlastfoot
KNN & 0.8389 & 0.8727 & 0.8389 & 0.8536 \\
RandomForest & 0.734 & 0.734 & 0.734 & 0.7047 \\
DecisionTree & 0.7603 & 0.7648 & 0.7603 & 0.7603 \\
CNN without Attention & 0.8984 & 0.9001 & 0.8984 & 0.8959 \\
CNN with Attention & \textbf{0.9066} & \textbf{0.907} & \textbf{0.9066}
& \textbf{0.9066} \\
\end{longtable}
}

\begin{quote}
\emph{\textbf{Bảng 3.4 Bảng so sánh kết quả giữa các giải thuật trên tập
dữ liệu Cifar-10}}
\end{quote}

{\def\LTcaptype{none} % do not increment counter
\begin{longtable}[]{@{}
  >{\raggedright\arraybackslash}p{(\linewidth - 8\tabcolsep) * \real{0.3863}}
  >{\raggedright\arraybackslash}p{(\linewidth - 8\tabcolsep) * \real{0.1339}}
  >{\raggedright\arraybackslash}p{(\linewidth - 8\tabcolsep) * \real{0.1337}}
  >{\raggedright\arraybackslash}p{(\linewidth - 8\tabcolsep) * \real{0.1489}}
  >{\raggedright\arraybackslash}p{(\linewidth - 8\tabcolsep) * \real{0.1181}}@{}}
\toprule\noalign{}
\begin{minipage}[b]{\linewidth}\centering
\textbf{Algorithm}
\end{minipage} & \begin{minipage}[b]{\linewidth}\raggedright
\begin{quote}
\textbf{Accuracy}
\end{quote}
\end{minipage} & \begin{minipage}[b]{\linewidth}\raggedright
\begin{quote}
\textbf{Precision}
\end{quote}
\end{minipage} & \begin{minipage}[b]{\linewidth}\raggedright
\begin{quote}
\textbf{Recall}
\end{quote}
\end{minipage} & \begin{minipage}[b]{\linewidth}\centering
\textbf{F1}
\end{minipage} \\
\midrule\noalign{}
\endhead
\bottomrule\noalign{}
\endlastfoot
KNN & 0.3567 & 0.4385 & 0.3567 & 0.3466 \\
RandomForest & 0.4356 & 0.4356 & 0.4356 & 0.4267 \\
DecisionTree & 0.311 & 0.3101 & 0.311 & 0.311 \\
CNN without Attention & 0.7157 & 0.7282 & 0.7157 & 0.714 \\
CNN with Attention & \textbf{0.7265} & \textbf{0.737} & \textbf{0.7265}
& \textbf{0.7292} \\
\end{longtable}
}

\begin{quote}
Nhìn vào bảng 3.3 và 3.4 có thể thấy được mô hình CNN có áp dụng cơ chế
attention CBAM mang lại hiệu xuất tốt hơn so với các mô hình không áp
dụng với các chỉ số như accuracy, precision, recall và f1 score. Trong
cả hai bảng mô hình CNN áp dụng cơ chế attention đều mang lại hiệu quả
tốt hơn so với các mô hình không áp dụng. Đặc biệt trong bảng 3.4 các
phương pháp học truyền thống như KNN, RandomForest và DecsionTree cho
thấy hiệu suất không ổn định và thấp, điều này chỉ ra rằng phương pháp
học truyền thống không đủ mạnh để giải quyết các bài toán phân loại hình
ảnh yêu cầu độ phức tạp cao. Ngược lại trong cả 2 bảng mô hình CNN đều
mang lại hiệu quả khá tốt, đặc biệt là khi áp dụng cơ chế attention, mô
hình CNN có áp dụng cơ chế attention đều mang lại accuracy, precision,
recall và f1 score cao hơn so với các phương pháp còn lại. Điều này cho
thấy cơ chế attention có tác động tích cực đến đến hiệu suất mô hình. Có
thể thấy được CBAM giúp mô hình tập trung vào các đặc trưng quan trọng
trong hình ảnh. Cơ chế Channel Attention làm nổi bật các

kênh chứa thông tin quan trọng, trong khi Spatial Attention giúp mô hình
tập trung vào các vùng không gian quan trọng trong hình ảnh. Sự kết hợp
này giúp cải thiện hiệu suất của mô hình một cách toàn diện. Qua đó có
thể chứng minh rằng học sâu kết hợp với cơ chế chú ý có thể cung cấp
hiệu suất phân loại tốt hơn đáng kể so với các phương pháp truyền thống.
\end{quote}

\paragraph{3.3.2.3. Kết quả thực nghiệm và biểu đồ Loss/Accuracy trên tập dữ liệu thực tiễn (PlantVillage, Chest-Xray-Pneumonia)}\label{3-3-2-3-practical-results}
\begin{quote}
\emph{[Nội dung phần này đang được cập nhật/viết tiếp...]}
\end{quote}

\paragraph{3.3.2.4. Phân tích hiện tượng suy giảm hiệu năng cục bộ (CBAM-Local) và ưu thế của giải pháp đề xuất (CBAM-Hook)}\label{3-3-2-4-performance-drop}
\begin{quote}
\emph{[Nội dung phần này đang được cập nhật/viết tiếp...]}
\end{quote}

\paragraph{3.3.2.5. Giải thích mô hình bằng trực quan hóa bản đồ nhiệt Grad-CAM (Sự tập trung vùng không gian đặc trưng thực tế)}\label{3-3-2-5-gradcam}
\begin{quote}
\emph{[Nội dung phần này đang được cập nhật/viết tiếp...]}
\end{quote}

\section{CHƯƠNG 4: KẾT QUẢ, KẾT LUẬN VÀ KIẾN NGHỊ}\label{chux1b0ux1a1ng-4-kux1ebft-quux1ea3-kux1ebft-luux1eadn-vuxe0-kiux1ebfn-nghux1ecb}

\subsection{4.1. KẾT QUẢ ĐẠT ĐƯỢC CỦA ĐỀ TÀI}\label{kux1ebft-quux1ea3}

\begin{quote}
Trong bài tiểu luận này, thông qua các tài liệu về mô hình mạng nơ-ron
tích chập (CNN) và cơ chế chú ý CBAM, học viên đã hiểu kiến trúc, cách
thức hoạt động của mô hình CNN, cùng với đó là cơ chế chú ý và áp dụng
vào mô hình CNN. Học viên đã khảo sát và hoàn thành thực hiện nghiên cứu
việc áp dụng cơ chế chú ý CBAM vào mô hình CNN với pre-trained model
VGG16 trên mặt lý thuyết lẫn thực nghiệm trên hai tập dữ liệu
MNIST-Fashion và Cifar-10. Học viên đã xây dựng được cơ chế chú ý CBAM
áp dụng vào sau một vài lớp của mô hình CNN và thực hiện huấn luyện cùng
với một số mô hình không áp dụng. Từ các số liệu từ hình 3.8,3.9 và bảng
3.3,3.4 so sánh accuracy, precision, recall và f1 score thu được thông
qua thực nghiệm đã cho thấy được mô hình CNN có áp dụng cơ chế chú ý
CBAM đã mang lại kết quả tích cực cho mô hình.
\end{quote}

\subsection{4.2. KẾT LUẬN (Đánh giá mức độ hoàn thành so với mục tiêu và ý nghĩa thực tiễn)}\label{kux1ebft-luux1eadn}

\begin{quote}
Qua những thông số thu được từ kết quả thực nghiệm cho thấy kết quả của
mô hình CNN thật sự tốt hơn so với các mô hình máy học truyền thống, đặc
biệt là việc áp dụng cơ chế chú ý CBAM vào mô hình đã cải thiện hiệu
suất cho mô hình thông qua các chỉ số như accuracy, precision, recall và
f1 score. Việc tích hợp CBAM vào mô hình CNN đã chứng minh được hiệu quả
trong việc nâng cao độ chính xác và hiệu suất tổng thể của mô hình phân
loại hình ảnh. Cơ chế Attention giúp mô hình tập trung vào các đặc trưng
và vùng quan trọng, cải thiện khả năng tổng quát hóa. Những kết quả này
khẳng định rằng việc sử dụng cơ chế chú ý à một hướng đi triển vọng
trong việc phát triển các mô hình học sâu cho bài toán phân loại hình
ảnh.
\end{quote}

\subsection{4.3. KIẾN NGHỊ VÀ CÁC HƯỚNG PHÁT TRIỂN TIẾP THEO}\label{kiux1ebfn-nghux1ecb}

\begin{quote}
Cơ chế chú ý đã thật sự mang lại hiệu quả tốt cho mô hình CNN, vì vậy
thông qua nghiên cứu này có thể nghiên cứu thêm các cơ chế chú ý khác và
áp dụng nhiều cơ chế chú ý khác nhau để tận dụng các ưu điểm của từng
loại để đạt được hiệu suất cao hơn. Nghiên cứu các phương pháp tối ưu
hóa để giảm thiểu chi phí tính toán của

cơ chế chú ý mà vẫn giữ được hiệu suất cao. Nghiên cứu mở rộng việc áp
dụng cơ chế chú ý vào các nhiệm vụ và các mô hình học sâu.
\end{quote}

\section{TÀI LIỆU THAM KHẢO}\label{tuxe0i-liux1ec7u-tham-khux1ea3o}

\begin{enumerate}
\def\labelenumi{\arabic{enumi}.}
\item
  Ashish Vaswani, Noam Shazeer, Niki Parmar, Jakob Uszkoreit, Llion
  Jones, Aidan N. Gomez, Łukasz Kaiser, Illia Polosukhin ``Attention is
  all you need'', 2017.
\item
  Trang web \ul{What is Convolutial Neural Network? Begin's Tutorial for
  Machine} \ul{Learning and Deep Learning}, truy cập ngày 12/03/2024.
\item
  Altenberger, Felix, and Claus Lenz. ``A non-technical survey on deep
  convolutional neural network architectures.'' arXiv preprint arXiv:
  1803.02129 (2018).
\item
  Cao, Changyu et al. ``A convolutional neural network face recognition
  algorithm based on data augmentation.'' (2019)
\item
  Simonyan, K.; Zisserman, A. Very deep convolutional networks for
  large-scale image recognition. arXiv 2014, arXiv:1409.1556..
\item
  Trang web \ul{Thuật toán CNN là gì? Và cấu trúc mạng Convolution
  Neural} \ul{Network}, truy cập ngày 5/4/2024.
\item
  Guo, M.-H., Xu, T.-X., Liu, J.-J., Liu, Z.-N., Jiang, P.-T., Mu,
  T.-J., ... \&
\end{enumerate}

\begin{quote}
Wang, Y. (2022). Attention mechanisms in computer vision: A survey.
Computational Visual Media, 8, 331-368.
\end{quote}

\begin{enumerate}
\def\labelenumi{\arabic{enumi}.}
\setcounter{enumi}{7}
\item
  Trang web \ul{Attention Mechanisms in Computer Vision: CBAM}, truy cập
  ngày 10/05/2024.
\item
  Sanghyun Woo, Jongchan Park, Joon-Young Lee, In So Kweon ``CBAM:
  Convolution Block Attention Module''.
\item
  Jie Hu, Li Shen, Samuel Albanie, Gang Sun, Enhua Wu
  ``Squeeze-and-Excitation Networks''.
\item
  Trang web \ul{CBAM-keras}, truy cập ngày 15/05/2024.
\item
  Trang web \ul{Understanding Transfer Learning for Deep Learning}, truy
  cập ngày 15/05/2024.
\item
  Trang web \ul{Fine tuning pre-trained model trong pytorch và áp dụng
  vào} \ul{Visual Saliency Prediction}, truy cập ngày 16/05/2024.
\end{enumerate}

\end{document}
